% concatenated from draft1_simulation_v1.tex
% \documentclass[final,times]{elsarticle}
% \documentclass{amsart}
\documentclass[preprint,times,3p]{elsarticle}

\usepackage{amsmath}
\usepackage{amssymb}
\usepackage{amsfonts}
\usepackage{graphicx}
\usepackage{epstopdf}
\usepackage{color}
\usepackage{bm}
\usepackage{multirow}
\usepackage{natbib}
\usepackage{hyperref}
\usepackage{cleveref}
\usepackage{amsthm}
% \usepackage{geometry}
\usepackage{lineno} 
\nolinenumbers

\usepackage{color}
\usepackage{float}
\usepackage{caption}

\usepackage[caption=false]{subfig}

\usepackage{ifpdf}
\ifpdf%
\usepackage{pdflscape}
\else
\usepackage{lscape}
\fi

\newcommand{\red}{\color{red}}

\hypersetup{
	colorlinks,
	linkcolor={red!50!black},
	citecolor={blue!50!black},
	urlcolor={blue!80!black}
}

%\usepackage{geometry}

% \geometry{a4paper,left=2.8cm,right=2.8cm,top=2.8cm,bottom=2.8cm}

% \theoremstyle{remark}  
\newtheorem{remark}{Remark} 
\newtheorem{lemma}{Lemma} 

\DeclareMathOperator{\Tr}{Tr}
\DeclareMathOperator{\tr}{tr}
\DeclareMathOperator{\spp}{sp}
\DeclareMathOperator{\liq}{li_2}
\DeclareMathOperator{\lif}{li_4}
\DeclareMathOperator{\lik}{li_k}
\DeclareMathOperator{\sgn}{sgn} 
\DeclareMathOperator{\Imag}{Im}
\DeclareMathOperator{\im}{Im}
\DeclareMathOperator{\Real}{Re} 
\DeclareMathOperator{\re}{Re} 

\def\A{\mbox{\boldmath $A$}}
\def\a{\mbox{\boldmath $a$}}
\def\B{\mbox{\boldmath $B$}}
\def\D{\mbox{\boldmath $D$}}
\def\E{\mbox{\boldmath $E$}}
\def\e{\mbox{\boldmath $e$}}
\def\f{\mbox{\boldmath $f$}}
\def\F{\mbox{\boldmath $F$}}
\def\g{\mbox{\boldmath $g$}}
\def\h{\mbox{\boldmath $h$}}
\def\m{\mbox{\boldmath $m$}}
\def\n{\mbox{\boldmath $n$}}
\def\M{\mbox{\boldmath $M$}}
\def\p{\mbox{\boldmath $p$}}
\def\q{\mbox{\boldmath $q$}}
\def\x{\mbox{\boldmath $x$}}
\def\y{\mbox{\boldmath $y$}}
\def\w{\mbox{\boldmath $w$}}
\def\bnu{\mbox{\boldmath $\nu$}}
\def\bell{\mbox{\boldmath $\ell$}}
\newcommand\dt {{\Delta t}}
\def\0{\mbox{\boldmath $0$}}
\def\um{\underline{\bm m}}

\def\ppsi{\mbox{\boldmath $\psi$}}


%\journal{Annals of Physics}
% \biboptions{sort&compress}



% \sloppy
\begin{document}

\begin{frontmatter}
% \title{An Energy Stable Structure Preserving Method for Micromagnetics Simulations}

\title{A New Fractional Step Structure Preserving Method for The Landau-Lifshitz-Gilbert Equation}

\author[XJTLU]{Changjian Xie\corref{cor}}
\cortext[cor]{Corresponding author.} 
\ead{Changjian.Xie@xjtlu.edu.cn}

% \author[XJTLU]{Yingxi Miao}
% \ead{Yingxi.Miao22@student.xjtlu.edu.cn}

% \author[XJTLU]{Haocheng Yang}
% \ead{Haocheng.Yang22@student.xjtlu.edu.cn}

\address[XJTLU]{School of Mathematics and Physics, Xi'an-Jiaotong-Liverpool University, Re'ai Rd. 111, Suzhou, 215123, Jiangsu, China.}



\begin{abstract}
In this paper, we propose a structure preserving method using a Crank-Nicolson's type method with an implicit Gauss-Seidel fractional iteration. Such a method is of first-order accuracy in time and second-order accuracy in space, stable and length preserving. Such kind of method brings great benefits for the theoretical analysis. The numerical accuracy, norm preserving and stability are verified for 1D and 3D tests. 

% One of the main difficulties in micromagnetics simulation is the severe time step
% constraint introduced by the exchange field. Using standard explicit integrators leads
% to a physical time step of sub-pico seconds, which is often two orders of magnitude
% smaller than the fastest physical time scales. Direct implicit integrators require solving complicated, coupled systems. Another major difficulty with the projection method in this field is the lack of rigorous theoretical guarantees regarding its stability of the projection step. In this paper, we introduce an implicit method
% whose complexity is comparable to solving the scalar heat equation implicitly. This
% method is structure preserving based on a combination of a Gauss–Seidel iteration and a Crank-Nicolson method. Such method allows us to carry out fully resolved calculations
% for the switching of the magnetization in micron-sized elements. 
\end{abstract}

\begin{keyword}
{Landau-Lifshitz-Gilbert equation\sep structure preserving method\sep energy stability}
\end{keyword}

\end{frontmatter}

\section{Introduction}


% Ferromagnetic materials find widespread application in data storage technologies, largely due to the bistable nature of their intrinsic magnetic ordering, commonly referred to as magnetization. 
% The temporal evolution of magnetization in these systems is principally dictated by 
The Landau-Lifshitz-Gilbert (LLG) equation \cite{Landau1935On,Gilbert:1955} incorporates two fundamental terms governing magnetization dynamics: a gyromagnetic term, responsible for energy conservation, and a damping term, which models energy dissipation. 
% The damping term is critically important in magnetic devices, profoundly influencing both their energy consumption and operational speed. Recent experimental work on magnetic-semiconductor heterostructures \cite{Zhang2020ExtremelyLM} has confirmed that the Gilbert damping constant can be tuned. At the microscopic scale, damping arises from various physical processes, including electron scattering, itinerant electron relaxation \cite{Heinrich1967TheIO}, and phonon-magnon coupling \cite{Suhl1998TheoryOT, Nan2020ElectricfieldCO}, which can be quantitatively assessed through electronic structure calculations \cite{TangXia2017}. From an applied standpoint, modulating the damping parameter enables optimization of magnetodynamic properties in ferromagnetic materials—for instance, by lowering the switching current and enhancing the writing speed of magnetic memory devices \cite{Wei2012MicromagneticsAR}. Although many experimental studies have focused on systems with small damping parameters \cite{Budhathoki2020LowGD,Lattery2018LowGD,Weber2019GilbertDO}, substantial damping effects have also been reported in works like \cite{GilbertKelly1955, Tanaka2014MicrowaveAssistedMR}. Notably, Tanaka et al. \cite{Tanaka2014MicrowaveAssistedMR} found that larger damping constants correlate with reduced magnetization switching times. Gilbert and Kelly \cite{GilbertKelly1955} also reported observing very large damping parameters (around 9) in their investigations. 
The LLG equation constitutes a vector-valued, nonlinear system characterized by a point-wise constant magnitude constraint on the magnetization vector. The model without damping is given by
\begin{align*}
    \m_t=-\m\times \Delta \m.
\end{align*}
Significant research efforts have been dedicated to devising efficient and numerically stable methods for micromagnetic simulations, as summarized in review articles such as \cite{kruzik2006recent,cimrak2007survey}. 
Such a equation can be solved with Crank-Nicolson's method (midpoint type) below,
\begin{align*}
    \frac{\m_h^{n+1}-\m_h^n}{\Delta t}=-\frac{\m_h^{n+1}+\m_h^n}{2}\times \Delta_h \left(\frac{\m_h^{n+1}+\m_h^n}{2}\right),
\end{align*}
which is a nonlinear scheme with great stability. However, this paper focus on the linearity, efficiency, stability and structure preservation. 
% Semi-implicit schemes have attracted considerable interest among existing numerical strategies due to their ability to maintain good numerical stability without requiring complex nonlinear solvers \cite{alouges2006convergence, gao2014optimal, Xie2018}. For instance, our research group previously proposed a second-order backward differentiation formula (BDF2) scheme based on one-sided interpolation \cite{Xie2018}. This method requires solving a three-dimensional linear system with non-constant coefficients at each time step. A rigorous theoretical analysis establishing the second-order convergence of this BDF2 method was provided in \cite{jingrun2019analysis}. Alternatively, Alouges et al. \cite{alouges2006convergence}  introduced a linearly implicit method utilizing the tangent space to preserve the magnetization constraint, although this approach achieves only first-order temporal accuracy. More recently, Lubich et al. \cite{Lubich2021} developed and analyzed high-order BDF schemes for the LLG equation. Unconditional unique solvability for such semi-implicit schemes has been established in \cite{jingrun2019analysis,Lubich2021}; however, their convergence analysis typically requires the temporal step size to be proportional to the spatial grid size. Despite these developments, the practical implementation and evaluation of high-order numerical methods for micromagnetic modeling continue to be active research areas. 
A notable limitation persists: the implicit Crank-Nicolson's method \cite{JEONG2010613} involves a nonlinear system with the norm preserving property. In this paper, we use such a method and such property to construct a structure preserving method based on implicit Gauss-Seidel iteration.

The rest of this paper is organized as follows. \cref{sec: numerical scheme} begins with a review of the micromagnetic model, followed by a detailed description of the proposed numerical scheme. \cref{sec:experiments} presents extensive numerical results, encompassing verification of temporal and spatial accuracy in one-dimensional (1D) and three-dimensional (3D) settings. 
% assessments of computational efficiency (via comparisons with BDF1 and BDF2), stability analysis with respect to the damping parameter, and an investigation of domain wall velocity dependence on both damping and external magnetic field. 
Concluding remarks and potential future research directions are provided in \cref{sec:conclusions}.


%\section{Formalism\label{Sec.Form}}
\section{The governing equation and numerical scheme}
\label{sec: numerical scheme}

%\subsection{Governing equation}

The Landau-Lifshitz-Gilbert (LLG) equation forms the fundamental basis of micromagnetics, providing a rigorous description of the spatiotemporal evolution of magnetization in ferromagnetic materials by incorporating two key physical phenomena: gyromagnetic precession and dissipative relaxation \cite{Landau1935On,Brown1963micromagnetics}. In nondimensional form, this governing equation is expressed as
% \begin{align}\label{c1-large}
% {\m}_t =-{\m}\times{\bm h}_{\text{eff}}-\alpha{\m}\times({\m}\times{\bm h}_{\text{eff}})
% \end{align}
\begin{align}\label{eq-5}
\m_t=-\m\times\Delta\m-\alpha\m\times(\m\times\Delta\m)
\end{align}
subject to the homogeneous Neumann boundary condition
\begin{equation}\label{boundary-large}
\frac{\partial{\m}}{\partial {\bm \nu}}\Big|_{\partial \Omega}=0,
\end{equation}
where \(\Omega \subset \mathbb{R}^d\) (\(d=1,2,3\)) represents the bounded domain of the ferromagnetic material, and \(\bm \nu\) is the unit outward normal vector on the boundary \(\partial \Omega\). This boundary condition ensures no magnetic surface charge, a physically appropriate assumption for isolated ferromagnetic systems.

The magnetization field \(\m: \Omega \to \mathbb{R}^3\) is a three-dimensional vector field satisfying the pointwise constraint \(|\m| = 1\),  stemming from the quantum mechanical alignment of electron spins in ferromagnets. The first term on the right-hand side of \cref{eq-5} describes gyromagnetic precession, where magnetic moments precess around the exchange field $\Delta \m$. The second term represents dissipative relaxation, with \(\alpha > 0\) being the dimensionless Gilbert damping coefficient that governs the rate of energy dissipation into the lattice.

% From the perspective of the Gibbs free energy functional, the effective field \(\bm h_{\text{eff}}\) is obtained as the functional derivative of the Gibbs free energy \(F[\m]\) with respect to the magnetization, i.e., \(\bm h_{\text{eff}} = -\delta F/\delta \m\). his functional incorporates all relevant energy contributions in ferromagnetic systems—exchange, anisotropy, magnetostatic (stray), and Zeeman energies—and is given by
% \begin{equation}\label{LL-Energy}
% F[\m] = \frac {\mu_0 M_s^2}{2} \left\{\int_\Omega \left( \epsilon|\nabla\m|^2 +
% q\left(m_2^2 + m_3^2\right)
% -2\h_e\cdot\m - \h_s\cdot\m \right)\mathrm{d}\x \right\} . 
% \end{equation}
% Here, \(\mu_0 = 4\pi \times 10^{-7}\, \text{H/m}\) is the vacuum permeability, \(M_s\) is the saturation magnetization, and \(\epsilon\) and \(q\) are dimensionless parameters defined subsequently. The vectors \(\h_e\) and \(\h_s\) denote the externally applied magnetic field and the stray field, respectively. For uniaxial ferromagnetic materials with a single easy axis, the effective field \(\bm h_{\text{eff}}\) decomposes into distinct physical components, yielding the explicit expression
% \begin{align}
% {\bm h}_{\text{eff}} =\epsilon\Delta\m-q(m_2\e_2+m_3\e_3)+\h_s+\h_e,
% \end{align}
% where \(\epsilon = C_{\text{ex}}/(\mu_0 M_s^2 L^2)\) and \(q = K_u/(\mu_0 M_s^2)\). Here, \(L\) is the characteristic length scale, \(C_{\text{ex}}\) is the exchange constant (controlling short-range spin alignment), and \(K_u\) is the uniaxial anisotropy constant (representing the energy cost for deviation from the easy axis). The unit vectors \(\e_2 = (0,1,0)\) and \(\e_3 = (0,0,1)\) efine the hard axes perpendicular to the uniaxial easy axis, and \(\Delta\) is the Laplacian operator in \(d\)-dimensional space.
% For Permalloy (NiFe), a common soft ferromagnetic material in spintronics, standard parameter values from the literature are: \(C_{\text{ex}} = 1.3 \times 10^{-11}\, \text{J/m}\), \(K_u = 100\, \text{J/m}^3\), and \(M_s = 8.0 \times 10^5\, \text{A/m}\). The stray field \(\h_s\) originates from magnetic charge distributions at domain boundaries and material surfaces and is mathematically represented by the integral expression
% \begin{align}\label{eqn:div}
% {\h}_{\text{s}}=\frac{1}{4\pi}\nabla \int_{\Omega} \nabla\left( \frac{1}{|\x-\y|}\right)\cdot {\bm m}({\bm y})\,d{\bm y},
% \end{align}
% which is a formulation that exhibits long-range spatial correlations. A critical computational advancement for practical micromagnetic simulations is that for rectangular domains \(\Omega\), the evaluation of \(\h_s\) can be efficiently computed via the Fast Fourier Transform (FFT) \cite{Wang2000}, which reduces the asymptotic computational complexity from \(O(N^d)\) to \(O(N^d \log N)\) for \(d\)-dimensional grids, enabling large-scale simulations.

% To facilitate numerical discretization, we introduce the composite source term
% \begin{align}\label{eq-4}
% \f=-Q(m_2\e_2+m_3\e_3)+\h_s+\h_e.
% \end{align}
% which aggregates the anisotropy, stray field, and external field contributions. Substituting this source term into \cref{c1-large}, the LLG equation is re-expressed as
% \begin{align}\label{eq-5}
% \m_t=-\m\times(\epsilon\Delta\m+\f)-\alpha\m\times(\m\times(\epsilon\Delta\m+\f).)
% \end{align}
% Leveraging the vector triple product identity \(\a \times ({\bm b} \times {\bm c}) = (\a \cdot {\bm c}){\bm b} - (\a \cdot {\bm b}){\bm c}\) and the pointwise constraint \(|\m| = 1\) (which implies \(\m \cdot \partial_t \m = 0\) via time differentiation), we simplify \cref{eq-5} to an equivalent formulation that is more amenable to stable numerical discretization:
% \begin{equation}\label{eq-model}
% \m_t=\alpha  (\epsilon\Delta\m+\f)+\alpha \left(\epsilon |\nabla \m|^2 -\m \cdot\f \right)\m-\m\times(\epsilon\Delta\m+\f).
% \end{equation}
%We establish a standardized discretization framework to underpin subsequent numerical approximations, defining temporal and spatial discretization notations and boundary condition enforcement strategies.

%Let \(k = t^{n+1} - t^n\) denote the uniform temporal step-size, with discrete time levels given by \(t^n = nk\) for \(n = 0,1,\dots,N_T\), where \(N_T = \lfloor T/k \rfloor\) and \(T\) is the final simulation time. For spatial discretization, a uniform Cartesian grid is employed with mesh-size \(h_x = h_y = h_z = h = L/N\) (for cubic domains of characteristic length \(L\)). The notation \(\m_{i,j,\ell}^n\) denotes the numerical approximation of \(\m(x_{i-1/2}, y_{j-1/2}, z_{\ell-1/2}, t^n)\), where \(x_{i-1/2} = (i - 1/2)h_x\), \(y_{j-1/2} = (j - 1/2)h_y\), and \(z_{\ell-1/2} = (\ell - 1/2)h_z\) define cell-centered grid positions. The index range \(-1 \leq i,j,\ell \leq N+2\) is adopted to accommodate boundary extrapolation.
%To enforce the homogeneous Neumann boundary condition \cref{boundary-large} while preserving high-order accuracy, a third-order extrapolation scheme is implemented. For instance, along the \(z\)-direction (normal to the boundary at \(z=0\) and \(z=L\)), the extrapolation rules for the magnetization field are:
%\begin{align*}
%\m_{i,j,1}=\m_{i,j,0},\quad \m_{i,j,-1}=\m_{i,j,2},\quad \m_{i,j,N+1}=\m_{i,j,N} ,\quad \m_{i,j,N+2}=\m_{i,j,N-1}.
%\end{align*}
%Analogous extrapolation formulas are applied to enforce the boundary condition along the \(x\)- and \(y\)-directions, ensuring consistent high-order accuracy across the entire computational domain.

%The fourth-order spatial difference operators should be introduced further. 
%To achieve fourth-order spatial accuracy, which is essential for resolving fine magnetic structures (e.g., domain walls with nanoscale width) and minimizing discretization-induced errors, we employ long-stencil finite difference operators for first and second partial derivatives. For the \(x\)-direction, the fourth-order accurate operators for \(\partial_x\) (denoted \({\mathcal D}_{x,(4)}^1\)) and \(\partial_x^2\) (denoted \({\mathcal D}_{x,(4)}^2\)) are defined as:
%\begin{align*}
%	{\mathcal D}_{x,(4)}^1 f_{i,j,k} &= \tilde{D}_x ( 1 - \frac{h_x^2}{6} D_x^2 ) f_{i,j,k} 
%	\\
%	&= 
%	\frac{  f_{i-2,j,k} - 8 f_{i-1,j,k}  + 8 f_{i+1,j,k} - f_{i+2,j,k} }{12 h_x} ,  
%	\\
%	{\mathcal D}_{x,(4)}^2 f_{i,j,k} &= D_x^2 ( 1 - \frac{h_x^2}{12} D_x^2 ) f_{i,j,k} 
%	\\
%	&=
%	\frac{ - f_{i-2,j,k} + 16 f_{i-1,j,k,k} - 30 f_{i,j,k} + 16 f_{i+1,j,k} - f_{i+2,j,k} }{12 h_x^2 } . 
%\end{align*}
%These operators are derived by canceling leading-order discretization errors via inclusion of adjacent grid points, resulting in fourth-order convergence (\(O(h^4)\)) for sufficiently smooth functions.
%
%By symmetry, the fourth-order difference operators for the \(y\)- and \(z\)-directions—designated \({\mathcal D}_{y,(4)}^1\), \({\mathcal D}_{y,(4)}^2\), \({\mathcal D}_{z,(4)}^1\), and \({\mathcal D}_{z,(4)}^2\)—are defined by substituting the respective spatial coordinates and mesh-sizes. The discrete Laplacian operator \(\Delta_h\), which approximates the continuous Laplacian \(\Delta\) with fourth-order accuracy, is constructed as the sum of the second-order operators in all three directions: \(\Delta_h = {\mathcal D}_{x,(4)}^2 + {\mathcal D}_{y,(4)}^2 + {\mathcal D}_{z,(4)}^2\).
%
%
%To provide a rigorous performance baseline for the proposed numerical method, we adapt the first-order (BDF1) and second-order (BDF2) backward differentiation formula schemes, which are standard in micromagnetics, to include the fourth-order spatial operators and the equivalent LLG formulation \cref{eq-model}. To better understand the process of stability, the proposed method is compared with an existing third order numerical scheme. These adapted BDF schemes will serve as benchmarks in subsequent numerical experiments, allowing quantitative evaluation of the proposed method's accuracy, efficiency, and stability.
 
% \section{Proposed method}

For the construction of the structure preserving method, we set $\alpha=0$ for \cref{eq-5}, we have the LLG equation below,
\begin{align}\label{eq-alpha-0}
\m_t=-\m\times\Delta\m.
\end{align}
It is obvious that \cref{eq-alpha-0} has a length-preserving property during the evolution process. To see this, we do scalar multiplication of \cref{eq-alpha-0} with $\m$:
\begin{align*}
    \frac{\partial \m}{\partial t} \cdot \m = -(\m \times \Delta \m) \cdot \m = 0.
\end{align*}
Then we have 
\begin{align*}
    \frac{\partial|\m|^2}{\partial t}=0,
\end{align*}
which implies $|\m(\x,t)|$ is constant for all $t$ and $\x$. And we assume that $|\m(\x,0)|=1$. Let 
Let \( E(\mathbf{m}(\mathbf{x}, t)) \) be an energy defined by \( E(\mathbf{m}(t)) := \|\nabla \mathbf{m}(t)\|_{L^2(\Omega)}^2 \). By taking an inner product of \cref{eq-alpha-0} with \( \Delta \mathbf{m} \), we obtain
\begin{align}\label{eq-3}
    \frac{\partial {\m}}{\partial t} \cdot \Delta {\m} = -({\m} \times \Delta {\m}) \cdot \Delta {\m} = 0. 
\end{align}
Using homogeneous Neumann or periodic boundary conditions, from \cref{eq-3} we have
\begin{align*}
    \begin{aligned}
0 &= \int_\Omega \frac{\partial {\m}}{\partial t} \cdot \Delta {\m}\, dx
   = \int_{\partial\Omega} \frac{\partial {\m}}{\partial t} \cdot \frac{\partial {\m}}{\partial {\n}}\, ds
   - \int_\Omega \nabla \frac{\partial {\m}}{\partial t} : \nabla {\m}\, dx \\
&= -\frac{1}{2} \frac{dE({\m}(t))}{dt},
\end{aligned}
\end{align*}
which implies that \(E({\m}(t))\) is constant and this problem has an energy conservation property. Here, \({\n}\) is a unit normal vector to \(\partial\Omega\) and the operator `:' is defined as \(A:B = \sum_{ij} a_{ij} b_{ij}\).
 
If we consider the simple linear vectorial equation 
\begin{align}\label{eq-CN}
\m_t=-\m\times \a,
\end{align}
where $\a^T=(a_1,a_2,a_3)$ is a constant vector. We use the Crank-Nicolson method to \cref{eq-CN}, we have
\begin{align}\label{eq-CN_1}
	\frac{\m_h^{n+1}-\m_h^n}{\Delta t}=-\frac{\m_h^{n+1}+\m_h^n}{2} \times \a,
\end{align}
which is norm preserving, since that if $\m_h^{n+1}+\m_h^n$ do the inner product for both sides, leading to
\begin{align*}
	\|\m_h^{n+1}\|_2=\|\m_h^n\|_2.
\end{align*}
After rearranging the compact form for \cref{eq-CN_1}, we have
\begin{align*}
	\begin{pmatrix}
	1&\frac12 \Delta t a_3&-\frac12 \Delta t a_2\\
	-\frac12 \Delta ta_3&1&\frac12 \Delta ta_1\\
	\frac12 \Delta ta_2&-\frac12 \Delta t a_1&1
	\end{pmatrix}\begin{pmatrix}
	m_1^{n+1}\\
	m_2^{n+1}\\
	m_3^{n+1}
	\end{pmatrix}=\begin{pmatrix}
	m_1^n+\frac12 \Delta t(a_2 m_3^n-a_3 m_2^n)\\
	m_2^n+\frac12 \Delta t(a_3 m_1^n-a_1 m_3^n)\\
	m_3^n+\frac12 \Delta t(a_1m_2^n-a_2 m_1^n)
	\end{pmatrix}
\end{align*}
We then have another form
\begin{align*}
	\begin{pmatrix}
	m_1^{n+1}\\
	m_2^{n+1}\\
	m_3^{n+1}
	\end{pmatrix}=\begin{pmatrix}
	1&\frac12 \Delta t a_3&-\frac12 \Delta t a_2\\
	-\frac12 \Delta ta_3&1&\frac12 \Delta ta_1\\
	\frac12 \Delta ta_2&-\frac12 \Delta t a_1&1
	\end{pmatrix}^{-1}\begin{pmatrix}
	1&-\frac12 \Delta t a_3 & \frac12 \Delta t a_3\\
	\frac12 \Delta t a_3&1& -\frac12 \Delta t a_1\\
	-\frac12 \Delta t a_2&\frac12 \Delta t a_1&1
	\end{pmatrix}\begin{pmatrix}
	m_1^n\\
	m_2^n\\
	m_3^n
	\end{pmatrix}=A\begin{pmatrix}
	m_1^n\\
	m_2^n\\
	m_3^n
	\end{pmatrix}
\end{align*}
where 
\begin{align*}
	A=\frac{1}{S}\begin{pmatrix}
	1+\beta^{2}a_1^{2}&-2\beta a_3+\beta^{2}a_1a_2&2\beta a_2+\beta^{2}a_1a_3\\
	2\beta a_3+\beta^{2}a_1a_2&1+\beta^{2}a_2^{2}&-2\beta a_1+\beta^{2}a_2a_3\\
	-2\beta a_2+\beta^{2}a_1a_3&2\beta a_1+\beta^{2}a_2a_3&1+\beta^{2}a_3^{2}
	\end{pmatrix}
\end{align*}
where $S=\det(A)=1+\beta^2(a_1^2+a_2^2+a_3^2)$ and $\beta=\frac{\Delta t}{2}$. The spectral analysis for the matrix $A$ is given by our previous work.

For \cref{eq-alpha-0}, we propose the following structure preserving schemes, which is a semi-implicit method:
	\begin{align}\label{eq-CN_4}
	\frac{\m_h^{n+1}-\m_h^n}{\Delta t}=-\frac{\m_h^{n+1}+\m_h^n}{2} \times \Delta_h \g_h,
	\end{align}
	where $\g_h^{s}=(I-\Delta t \Delta_h)^{-1}\m_h^s,\; s=n,n+1$. To be specific, we propose
	\begin{align*}
	\begin{pmatrix}
	1&\frac12 \Delta t \Delta_hg_3&-\frac12 \Delta t \Delta_hg_2\\
	-\frac12 \Delta t\Delta_hg_3&1&\frac12 \Delta t\Delta_hg_1\\
	\frac12 \Delta t\Delta_hg_2&-\frac12 \Delta t\Delta_h g_1&1
	\end{pmatrix}\begin{pmatrix}
	m_1^{n+1}\\
	m_2^{n+1}\\
	m_3^{n+1}
	\end{pmatrix}=\begin{pmatrix}
	m_1^n+\frac12 \Delta t(\Delta_hg_2 m_3^n-\Delta_hg_3 m_2^n)\\
	m_2^n+\frac12 \Delta t(\Delta_hg_3 m_1^n-\Delta_hg_1 m_3^n)\\
	m_3^n+\frac12 \Delta t(\Delta_hg_1m_2^n-\Delta_hg_2 m_1^n)
	\end{pmatrix}
	\end{align*}
	The trick is to handling the $\g$ evaluate at $t_n$ or $t_{n+1}$. If $\g$ evaluate at $t_n$, we have
	\begin{align*}
	\begin{pmatrix}
	1&\frac12 \Delta t \Delta_hg_3^n&-\frac12 \Delta t \Delta_hg_2^n\\
	-\frac12 \Delta t\Delta_hg_3^n&1&\frac12 \Delta t\Delta_hg_1^n\\
	\frac12 \Delta t\Delta_hg_2^n&-\frac12 \Delta t \Delta_hg_1^n&1
	\end{pmatrix}\begin{pmatrix}
	m_1^{n+1}\\
	m_2^{n+1}\\
	m_3^{n+1}
	\end{pmatrix}=\begin{pmatrix}
	m_1^n+\frac12 \Delta t(\Delta_hg_2^n m_3^n-\Delta_hg_3^n m_2^n)\\
	m_2^n+\frac12 \Delta t(\Delta_hg_3^n m_1^n-\Delta_hg_1^n m_3^n)\\
	m_3^n+\frac12 \Delta t(\Delta_hg_1^n m_2^n-\Delta_hg_2^n m_1^n)
	\end{pmatrix}
	\end{align*}
	which is proved to be mildly better than CFL condition. We choose the 1D exact solution as an example, and get the results presented in \Cref{tab-a-E-4}.
    \begin{table}[htbp]
	\centering
	\caption{The explicit scheme when $h = 5D-4$, $T=1d-1$ in 1D.}\label{tab-a-E-4}
	\begin{tabular}{|c|c|c|c|c|}
		\hline
		$k$ & $\|\m_h-\m_e\|_\infty$ & $\|\m_h-\m_e\|_2$ &$\|\m_h-\m_e\|_{H^1}$ & CPU time (s)\\
		\hline
		2.0D-2 & 0.018000157890128 &0.013706389798827&0.109658855967163&0.066217 \\
		1.0D-2 &0.012210502256479&0.008808376100206&0.069865327666244&0.115753 \\
		5.0D-3 &0.006883104572874&0.004488186175718&0.033485647423559&0.195630\\
		2.5D-3 &0.003297684177562&0.002226171926094&0.016241611196008&0.317750\\
		1.25D-3 &0.006813614081391&0.002600274876724&3.155993419074890&0.733073\\
		6.25D-4 &1.908302419811126 &1.414146167531504&1.379341310610379e+03&1.675463\\
		\hline
	\end{tabular}
\end{table}

    
    We propose a structure preserving method below:
    \begin{itemize}
        \item step 1. Solving the first linear system below,
        \begin{align*}
	\begin{pmatrix}
	1&\frac12 \Delta t \Delta_hg_3^n&-\frac12 \Delta t \Delta_hg_2^n\\
	0&1&0\\
	0&0&1
	\end{pmatrix}\begin{pmatrix}
	m_1^{n+1}\\
	m_2^{n+1}\\
	m_3^{n+1}
	\end{pmatrix}=\begin{pmatrix}
	m_1^n+\frac12 \Delta t(\Delta_hg_2^n m_3^n-\Delta_hg_3^n m_2^n)\\
	m_2^n+\frac12 \Delta t(\Delta_hg_3^n m_1^n-\Delta_hg_1^n m_3^n)\\
	m_3^n+\frac12 \Delta t(\Delta_hg_1^n m_2^n-\Delta_hg_2^n m_1^n)
	\end{pmatrix}
	\end{align*}
    \item step 2. Solving the second linear system below,
	\begin{align*}
	\begin{pmatrix}
	1&\frac12 \Delta t \Delta_hg_3^n&-\frac12 \Delta t \Delta_hg_2^n\\
	-\frac12 \Delta t\Delta_hg_3^n&1&\frac12 \Delta t\Delta_hg_1^{n+1}\\
	0&0&1
	\end{pmatrix}\begin{pmatrix}
	m_1^{n+1}\\
	m_2^{n+1}\\
	m_3^{n+1}
	\end{pmatrix}=\begin{pmatrix}
	m_1^n+\frac12 \Delta t(\Delta_hg_2^n m_3^n-\Delta_hg_3^n m_2^n)\\
	m_2^n+\frac12 \Delta t(\Delta_hg_3^n m_1^n-\Delta_hg_1^{n+1} m_3^n)\\
	m_3^n+\frac12 \Delta t(\Delta_hg_1^{n+1}m_2^n-\Delta_hg_2^n m_1^n)
	\end{pmatrix}
	\end{align*}
     \item step 3. Solving the third linear system below,
	\begin{align}\label{eq-step-3}
	\begin{pmatrix}
	1&\frac12 \Delta t \Delta_hg_3^n&-\frac12 \Delta t \Delta_hg_2^{n+1}\\
	-\frac12 \Delta t\Delta_hg_3^n&1&\frac12 \Delta t\Delta_hg_1^{n+1}\\
	\frac12 \Delta t\Delta_hg_2^{n+1}&-\frac12 \Delta t\Delta_h g_1^{n+1}&1
	\end{pmatrix}\begin{pmatrix}
	m_1^{n+1}\\
	m_2^{n+1}\\
	m_3^{n+1}
	\end{pmatrix}=\begin{pmatrix}
	m_1^n+\frac12 \Delta t(\Delta_hg_2^{n+1} m_3^n-\Delta_hg_3^n m_2^n)\\
	m_2^n+\frac12 \Delta t(\Delta_hg_3^n m_1^n-\Delta_hg_1^{n+1} m_3^n)\\
	m_3^n+\frac12 \Delta t(\Delta_hg_1^{n+1}m_2^n-\Delta_hg_2^{n+1} m_1^n)
	\end{pmatrix}
	\end{align}
    which is proved to be unconditionally stable and norm preserving with $\|\m_h^{n+1}\|_2=\|\m_h^n\|_2$.
    \end{itemize}

\begin{remark}
If we consider the explicit treatment for $\Delta \m$, we get 
	\begin{align}\label{eq-CN_2}
	\frac{\m_h^{n+1}-\m_h^n}{\Delta t}=-\frac{\m_h^{n+1}+\m_h^n}{2} \times \Delta_h \m_h^n,
	\end{align}
	which gives a CFL-type condition for the stability.
If we consider the implicit treatment for $\Delta \m$, we get
	\begin{align}\label{eq-CN_3}
	\frac{\m_h^{n+1}-\m_h^n}{\Delta t}=-\frac{\m_h^{n+1}+\m_h^n}{2} \times \Delta_h \m_h^{n+1},
	\end{align}
	which poses a difficulty for the nonlinear systems with high complexity.
\end{remark}

The proposed method for \cref{eq-5} with damping term is based on the semi-implicit method below,
	\begin{align}\label{eq-CN_4}
	\frac{\m_h^{n+1}-\m_h^n}{\Delta t}=-\frac{\m_h^{n+1}+\m_h^n}{2} \times \Delta_h \g_h-\alpha\frac{\m_h^{n+1}+\m_h^n}{2} \times (\m_h^n\times \Delta_h \g_h),
	\end{align}
	where $\g_h^{s}=(I-\Delta t \Delta_h)^{-1}\m_h^s,\; s=n,n+1$. 
We propose a structure preserving method below:
    \begin{itemize}
        \item step 1. Solving the first linear system below,
        \begin{align*}
	\begin{pmatrix}
	1&\frac12 \Delta t h_3^n&-\frac12 \Delta t h_2^n\\
	0&1&0\\
	0&0&1
	\end{pmatrix}\begin{pmatrix}
	m_1^{n+1}\\
	m_2^{n+1}\\
	m_3^{n+1}
	\end{pmatrix}=\begin{pmatrix}
	b_1^n\\
	b_2^n\\
	b_3^n
	\end{pmatrix}
	\end{align*}
    where 
    \begin{align*}
% h_1^n&=\Delta_h g_1^{n}+\alpha \left[m_2^n \cdot \Delta_hg_3^{n}-m_3^n\cdot \Delta_h g_2^{n}\right]\\
h_2^n&=\Delta_h g_2^{n}+\alpha \left[m_3^n \cdot \Delta_hg_1^{n}-m_1^n\cdot \Delta_h g_3^{n}\right],\\
h_3^n&=\Delta_h g_3^{n}+\alpha \left[m_1^n \cdot \Delta_hg_2^{n}-m_2^n\cdot \Delta_h g_1^{n}\right]
\end{align*}
and 
\begin{align*}
b_1^n&=m_1^n+\frac12 k [( \Delta_h g_2^{n})m_3^n-( \Delta_h g_3^{n})m_2^n]\\
&\quad+\frac12 \alpha k [(m_3^n(\Delta_h g_1^{n})-m_1^n(\Delta_h g_3^{n}))m_3^n-(m_1^n(\Delta_h g_2^{n})-m_2^n(\Delta_h g_1^{n}))m_2^n]\\
b_2^n&=m_2^n+\frac12 k [(\Delta_h g_3^{n})m_1^n-(\Delta_h g_1^{n})m_3^n]\\
&\quad+\frac12 \alpha k [(m_1^n(\Delta_h g_2^{n})-m_2^n(\Delta_h g_1^{n}))m_1^n-(m_2^n(\Delta_h g_3^{n})-m_3^n(\Delta_h g_2^{n}))m_3^n]\\
b_3^n&=m_3^n+\frac12 k [(\Delta_h g_1^{n})m_2^n-(\Delta_h g_2^{n})m_1^n]\\
&\quad+\frac12 \alpha k [(m_2^n(\Delta_h g_3^{n})-m_3^n(\Delta_h g_2^{n}))m_2^n-(m_3^n(\Delta_h g_1^{n})-m_1^n(\Delta_h g_3^{n}))m_1^n].
\end{align*}
    \item step 2. Solving the second linear system below,
	\begin{align*}
	\begin{pmatrix}
	1&\frac12 \Delta t h_3^n&-\frac12 \Delta t h_2^n\\
	-\frac12 \Delta th_3^n&1&\frac12 \Delta th_1^{n+1}\\
	0&0&1
	\end{pmatrix}\begin{pmatrix}
	m_1^{n+1}\\
	m_2^{n+1}\\
	m_3^{n+1}
	\end{pmatrix}=\begin{pmatrix}
	b_1^n\\
	b_2^n\\
	b_3^n
	\end{pmatrix}
	\end{align*}
        where 
    \begin{align*}
h_1^{n+1}&=\Delta_h g_1^{n+1}+\alpha \left[m_2^n \cdot \Delta_hg_3^{n}-m_3^n\cdot \Delta_h g_2^{n}\right]\\
h_2^n&=\Delta_h g_2^{n}+\alpha \left[m_3^n \cdot \Delta_hg_1^{n}-m_1^n\cdot \Delta_h g_3^{n}\right],\\
	h_3^n&=\Delta_h g_3^{n}+\alpha \left[m_1^n \cdot \Delta_hg_2^{n}-m_2^n\cdot \Delta_h g_1^{n}\right]
\end{align*}
and 
\begin{align*}
b_1^n&=m_1^n+\frac12 k [( \Delta_h g_2^{n})m_3^n-( \Delta_h g_3^{n})m_2^n]\\
&\quad+\frac12 \alpha k [(m_3^n(\Delta_h g_1^{n+1})-m_1^n(\Delta_h g_3^{n}))m_3^n-(m_1^n(\Delta_h g_2^{n})-m_2^n(\Delta_h g_1^{n+1}))m_2^n]\\
b_2^n&=m_2^n+\frac12 k [(\Delta_h g_3^{n})m_1^n-(\Delta_h g_1^{n+1})m_3^n]\\
&\quad+\frac12 \alpha k [(m_1^n(\Delta_h g_2^{n})-m_2^n(\Delta_h g_1^{n+1}))m_1^n-(m_2^n(\Delta_h g_3^{n})-m_3^n(\Delta_h g_2^{n}))m_3^n]\\
b_3^n&=m_3^n+\frac12 k [(\Delta_h g_1^{n+1})m_2^n-(\Delta_h g_2^{n})m_1^n]\\
&\quad+\frac12 \alpha k [(m_2^n(\Delta_h g_3^{n})-m_3^n(\Delta_h g_2^{n}))m_2^n-(m_3^n(\Delta_h g_1^{n+1})-m_1^n(\Delta_h g_3^{n}))m_1^n].
\end{align*}
     \item step 3. Solving the third linear system below,
	\begin{align*}
	\begin{pmatrix}
	1&\frac12 \Delta t h_3^n&-\frac12 \Delta t h_2^{n+1}\\
	-\frac12 \Delta th_3^n&1&\frac12 \Delta th_1^{n+1}\\
	\frac12 \Delta th_2^{n+1}&-\frac12 \Delta th_1^{n+1}&1
	\end{pmatrix}\begin{pmatrix}
	m_1^{n+1}\\
	m_2^{n+1}\\
	m_3^{n+1}
	\end{pmatrix}=\begin{pmatrix}
	b_1^n\\
	b_2^n\\
	b_3^n
	\end{pmatrix}
	\end{align*}
            where 
    \begin{align*}
h_1^{n+1}&=\Delta_h g_1^{n+1}+\alpha \left[m_2^n \cdot \Delta_hg_3^{n}-m_3^n\cdot \Delta_h g_2^{n+1}\right]\\
h_2^{n+1}&=\Delta_h g_2^{n+1}+\alpha \left[m_3^n \cdot \Delta_hg_1^{n+1}-m_1^n\cdot \Delta_h g_3^{n}\right],\\
	h_3^n&=\Delta_h g_3^{n}+\alpha \left[m_1^n \cdot \Delta_hg_2^{n+1}-m_2^n\cdot \Delta_h g_1^{n+1}\right]
\end{align*}
and 
\begin{align*}
b_1^n&=m_1^n+\frac12 k [( \Delta_h g_2^{n+1})m_3^n-( \Delta_h g_3^{n})m_2^n]\\
&\quad+\frac12 \alpha k [(m_3^n(\Delta_h g_1^{n+1})-m_1^n(\Delta_h g_3^{n}))m_3^n-(m_1^n(\Delta_h g_2^{n+1})-m_2^n(\Delta_h g_1^{n+1}))m_2^n]\\
b_2^n&=m_2^n+\frac12 k [(\Delta_h g_3^{n})m_1^n-(\Delta_h g_1^{n+1})m_3^n]\\
&\quad+\frac12 \alpha k [(m_1^n(\Delta_h g_2^{n+1})-m_2^n(\Delta_h g_1^{n+1}))m_1^n-(m_2^n(\Delta_h g_3^{n})-m_3^n(\Delta_h g_2^{n+1}))m_3^n]\\
b_3^n&=m_3^n+\frac12 k [(\Delta_h g_1^{n+1})m_2^n-(\Delta_h g_2^{n+1})m_1^n]\\
&\quad+\frac12 \alpha k [(m_2^n(\Delta_h g_3^{n})-m_3^n(\Delta_h g_2^{n+1}))m_2^n-(m_3^n(\Delta_h g_1^{n+1})-m_1^n(\Delta_h g_3^{n}))m_1^n].
\end{align*}
    which is proved to be unconditionally stable and norm preserving with $\|\m_h^{n+1}\|_2=\|\m_h^n\|_2$.
    \end{itemize}

% \section{Theoretical results}

% For the simplicity of the analysis, we take $\alpha=0$ and get the main results below.

% \subsection{stability}

% Consider the upper triangular matrix \( A \) for the first step of the three-stage linear system:
% \[
% A = \begin{pmatrix}
% 1 & \frac{1}{2}\Delta t \Delta_h g_3^n & -\frac{1}{2}\Delta t \Delta_h g_2^n \\
% 0 & 1 & 0 \\
% 0 & 0 & 1
% \end{pmatrix}
% \]
% A key property of upper triangular matrices is that their eigenvalues are precisely the elements on the main diagonal. For matrix \( A \), the diagonal entries are all \( 1 \), so the eigenvalues satisfy \( \lambda_1 = \lambda_2 = \lambda_3 = 1 \).

% The spectral radius \( \rho(A) \) of a matrix is defined as:
% \[
% \rho(A) = \max\left\{|\lambda_1|, |\lambda_2|, |\lambda_3|\right\}
% \]
% Substituting the eigenvalues of \( A \) into the definition, we get:
% \[
% \rho(A) = \max\left\{|1|, |1|, |1|\right\} = 1
% \]
% This result (\( \rho(A) = 1 \)) ensures the unconditional stability of the numerical scheme, as the solution remains bounded without exponential growth during iterative time stepping.

% \begin{align*}
%     \begin{pmatrix}
%         I & A\\
%         -A & I
%     \end{pmatrix}
% \end{align*}
% where $A$ is a diagonal matrix with different entries. $I$ is unit matrix. How do we get the spectral radius?

% What about CN scheme?

% \subsection{convergence rate}

\section{Scheme Property}
For simplicity, we analyze the case without damping. Combine with \cref{eq-CN_4} and \cref{eq-step-3}, we update satbilized $g_2^{n+1}$ and $g_3^{n+1}$ by first two steps. We still can use the \cref{eq-CN_4} and take the inner product with $\m_h^{n+1}+\m_h^n$. Hence we have
\begin{align*}
    |\m_h^{n+1}|^2=|\m_h^{n}|^2.
\end{align*}

\begin{remark}
    Notice that the Crank-Nicolson's scheme below,
\begin{align}\label{eq-6}
    \frac{{\m}^{n+1} - {\m}^n}{\Delta t} = -\frac{\m_h^{n+1}+\m_h^n}{2} \times \Delta_h \frac{\m_h^{n+1}+\m_h^n}{2} .
\end{align}
We can easily verify the scheme conserves the magnitude of magnetization. On the other hand, forming an inner product between \cref{eq-6} and $\Delta_h(\m_h^{n+1}+\m_h^n)$, we have
\begin{align}\label{eq-8}
    \frac{{\m}^{n+1} - {\m}^n}{\Delta t} \cdot \Delta_h (\m_h^{n+1}+\m_h^n) = -\left(\frac{\m_h^{n+1}+\m_h^n}{2} \times \Delta_h \frac{\m_h^{n+1}+\m_h^n}{2} \right)\cdot \Delta_h (\m_h^{n+1}+\m_h^n)=0 .
\end{align}
Summation \cref{eq-8} over $i=1,\cdots,N_x$, we have
\begin{align*}
   0=\sum_{i=1}^{N_x}\frac{{\m}_i^{n+1} - {\m}_i^n}{\Delta t} \cdot \frac{1}{h_x^2}\left(\m_{i+1}^{n+1}+\m_{i}^n-2(\m_i^{n+1}+\m_i^n)+\m_{i-1}^{n+1}+\m_{i-1}^n\right)
\end{align*}
where the periodic boundary condition is applied.
% where the homogeneous Neumann boundary condition is applied. Say $\m_0=\m_{1}$ and $\m_{N_x}=\m_{N_x+1}$. 
Now we get 
\begin{align*}
    \sum_{i=1}^{N_x} \left( 2{\m}_i^{n+1} \cdot {\m}_{i+1}^{n+1} - 2{\m}_i^n \cdot {\m}_{i+1}^n \right) = 0.
\end{align*}
Now, we get the following energy conservation:
\begin{align*}
    \begin{aligned}
E({\m}^{n+1}) - E({\m}^n) 
&= \frac{1}{h} \sum_{i=1}^{N_x} \left( |{\m}_{i+1}^{n+1} - {\m}_i^{n+1}|^2 - |{\m}_{i+1}^n - {\m}_i^n|^2 \right) \\
&= \frac{1}{h} \sum_{i=1}^{N_x} \left( |{\m}_{i+1}^{n+1}|^2 - |{\m}_{i+1}^n|^2 + |{\m}_i^{n+1}|^2 - |{\m}_i^n|^2 - 2{\m}_i^{n+1} \cdot {\m}_{i+1}^{n+1} + 2{\m}_i^n \cdot {\m}_{i+1}^n \right) = 0,
\end{aligned}
\end{align*}
where $|\m|=1$ has been used. 

The Crank-Nicolson method is sometimes introduced with the midpoint-type formulation \cref{eq-8} and the trapezoidal approach below,
\begin{align}\label{eq-8-tr}
    \frac{\m^{n+1}-\m^n}{\Delta t}=\frac12 \left[ (\m\times \Delta_h \m)^{n+1}+(\m\times \Delta_h \m)^n\right].
\end{align}
For the CN scheme \cref{eq-8-tr}, we can write as below,
\begin{align}\label{eq-11}
    \m^{n+1}+\frac{\Delta t}{2}(\m^{n+1}\times \Delta_h \m^{n+1})=\m^n-\frac{\Delta t}{2}(\m^n\times \Delta_h \m^n),
\end{align}
which can be rewrite as a component form,
\begin{align}\label{eq-13}
    \begin{pmatrix}
        u_i^{n+1}\\
        v_i^{n+1}\\
        w_i^{n+1}
    \end{pmatrix}+\frac{\Delta t}{2}\begin{pmatrix}
        v_i\Delta_h w_i-w_i\Delta_h v_i\\
        w_i\Delta_h u_i-u_i\Delta_h w_i\\
        u_i\Delta_h v_i-v_i\Delta_h u_i
    \end{pmatrix}^{n+1}=b_i^n,\quad \text{ for } i=1,2,\cdots,N_x,
\end{align}
where ${\bm b}=\m^n-\frac{\Delta t}{2}(\m^n\times \Delta_h \m^n)$. Note that
\begin{align}
v_i \Delta_h w_i - w_i \Delta_h v_i 
&= v_i \frac{w_{i-1} - 2w_i + w_{i+1}}{h_x^2} - w_i \frac{v_{i-1} - 2v_i + v_{i+1}}{h_x^2} \notag \\
&= v_i \frac{w_{i-1} + w_{i+1}}{h_x^2} - w_i \frac{v_{i-1} + v_{i+1}}{h_x^2} \notag \\
&= v_i \tilde{\Delta}_h w_i - w_i \tilde{\Delta}_h v_i, \quad \text{for } i = 1, \dots, N_x,
\label{eq:cancel1}
\end{align}
where \(\tilde{\Delta}_h w_i = (w_{i-1} + w_{i+1})/h_x^2\). Similarly, we have
\begin{align}
w_i \Delta_h u_i - u_i \Delta_h w_i &= w_i \tilde{\Delta}_h u_i - u_i \tilde{\Delta}_h w_i,
\label{eq:cancel2} \\
u_i \Delta_h v_i - v_i \Delta_h u_i &= u_i \tilde{\Delta}_h v_i - v_i \tilde{\Delta}_h u_i.
\label{eq:cancel3}
\end{align}

This cancelation stabilizes the scheme. By Eqs.~\eqref{eq:cancel1}, \eqref{eq:cancel2} and \eqref{eq:cancel3} we rewrite the above equation:
\begin{align}
A_i \begin{pmatrix}
u_i^{n+1} \\
v_i^{n+1} \\
w_i^{n+1}
\end{pmatrix}
= \begin{pmatrix}
\alpha_i \\
\beta_i \\
\gamma_i
\end{pmatrix},
\end{align}
where
\begin{align}
A_i = 
\begin{pmatrix}
1 & \frac{\Delta t}{2} \tilde{\Delta}_h w_i^{n+1} & -\frac{\Delta t}{2} \tilde{\Delta}_h v_i^{n+1} \\
-\frac{\Delta t}{2} \tilde{\Delta}_h w_i^{n+1} & 1 & \frac{\Delta t}{2} \tilde{\Delta}_h u_i^{n+1} \\
\frac{\Delta t}{2} \tilde{\Delta}_h v_i^{n+1} & -\frac{\Delta t}{2} \tilde{\Delta}_h u_i^{n+1} & 1
\end{pmatrix}
=
\begin{pmatrix}
1 & c & -b \\
-c & 1 & a \\
b & -a & 1
\end{pmatrix}
\end{align}
and \((\alpha_i, \beta_i, \gamma_i)^T\) is the right-hand side term in \cref{eq-13}. Then using Cramer's rule, we obtain
\begin{align}
(u_i^{n+1}, v_i^{n+1}, w_i^{n+1}) = 1/|A_i| \bigl(|A_{i,1}|, |A_{i,2}|, |A_{i,3}|\bigr), \quad i = 1, \dots, N_x,
\end{align}
where \(A_{i,j}\) is obtained by replacing the \(j\)-th column of \(A_i\) with \((\alpha_i, \beta_i, \gamma_i)^T\).
\begin{align*}
|A_i| &= 1 + a^2 + b^2 + c^2, \\
|A_{i,1}| &= \alpha_i(1 + a^2) - \beta_i(c - ab) + \gamma_i(ac + b), \\
|A_{i,2}| &= \alpha_i(ab + c) + \beta_i(1 + b^2) - \gamma_i(a - bc), \\
|A_{i,3}| &= \alpha_i(ac - b) + \beta_i(a + bc) + \gamma_i(1 + c^2).
\end{align*}

We can rewrite \cref{eq-11} as a matrix form:
\[
\begin{pmatrix}
1 & c & -b \\
-c & 1 & a \\
b & -a & 1
\end{pmatrix}
\mathbf{m}_i^{n+1} = \boldsymbol{\phi}_i^n,
\]
where
\[
a = \frac{\Delta t}{2} \tilde{\Delta}_h u_i^{n+1}, \quad
b = \frac{\Delta t}{2} \tilde{\Delta}_h v_i^{n+1}, \quad \text{and} \quad
c = \frac{\Delta t}{2} \tilde{\Delta}_h w_i^{n+1}.
\]

Indeed, when we consider the scheme without cancellation. We rewrite \cref{eq-13}.
\[
A_i
\begin{pmatrix}
u_i^{n+1} \\
v_i^{n+1} \\
w_i^{n+1}
\end{pmatrix}
= \boldsymbol{\phi}_i^n
\quad \text{for } i = 1, \dots, N_x.
\]

To discuss the stability of the Crank--Nicolson scheme, we compute the characteristic polynomial of \(A_i\).
\[
\det(A_i - \lambda I) = (1 - \lambda)^3 + (1 - \lambda)(a^2 + b^2 + c^2).
\]

The three eigenvalues of \(A_i\) are
\[
\lambda_1 = 1, \quad
\lambda_2 = 1 + \mathrm{i}\sqrt{a^2 + b^2 + c^2}, \quad
\text{and} \quad
\lambda_3 = 1 - \mathrm{i}\sqrt{a^2 + b^2 + c^2}.
\]
Thus, the three eigenvalues of \(A_i^{-1}\) are
\[
\gamma_1 = 1/\lambda_1, \quad \gamma_2 = 1/\lambda_2, \quad \text{and} \quad \gamma_3 = 1/\lambda_3.
\]

The absolute values of three eigenvalues of \(A_i^{-1}\) are
\[
|\gamma_1| = 1, \quad |\gamma_2| = |\gamma_3| = \frac{1}{\sqrt{1 + a^2 + b^2 + c^2}} \leq 1.
\]

Without cancellation, \(a\), \(b\), and \(c\) are small compared to 1, on the other hand, with cancelation \(a^2 + b^2 + c^2 \approx \mathcal{O}(1/h_x^4)\).
Therefore, \(1/\sqrt{1 + a^2 + b^2 + c^2} \approx \mathcal{O}(h_x^2) \ll 1\) and this makes the iterations stable. In short, the CN scheme is stable.

For the truncation error of \cref{eq-8}, Finally, we show that the truncation error of the scheme is second order in time and space. Let \(u, v, w\) be the exact solution of the LL equation without damping term and a forcing term. Then, the local truncation error of the first component of the equations is
\begin{align*}
\tau_i^{n+\frac{1}{2}}
&= \frac{u_i^{n+1} - u_i^n}{\Delta t}
+ \frac{v_i^{n+1} + v_i^n}{2} \Delta_h\left( \frac{w_i^{n+1} + w_i^n}{2} \right)
- \frac{w_i^{n+1} + w_i^n}{2} \Delta_h\left( \frac{v_i^{n+1} + v_i^n}{2} \right) \\
&= (u_t)_i^{n+\frac{1}{2}} + O(\Delta t^2)
+ \left( v_i^{n+\frac{1}{2}} + O(\Delta t^2) \right)\left( (w_{xx})_i^{n+\frac{1}{2}} + O(h^2) + O(\Delta t^2) \right) \\
&\quad - \left( w_i^{n+\frac{1}{2}} + O(\Delta t^2) \right)\left( (v_{xx})_i^{n+\frac{1}{2}} + O(h^2) + O(\Delta t^2) \right) \\
&= \left( u_t + v w_{xx} - w v_{xx} \right)_i^{n+\frac{1}{2}} + O(h^2) + O(\Delta t^2).
\end{align*}

Since \(u, v, w\) is the solution of the differential equation so
\[
\left( u_t + v w_{xx} - w v_{xx} \right)_i^{n+\frac{1}{2}} = 0.
\]

Therefore, the principal part of the local truncation error is
\[
\tau_i^{n+\frac{1}{2}} = O(h^2) + O(\Delta t^2).
\]

For the second and third components of the equations, we get same results.

For the stability analysis of \cref{eq-8-tr}, we take $\e^n=\m^n-\m(t^n)$. Substituting the exact solution \(\boldsymbol{m}(t^{n+1})\) into the scheme, the truncation error \(\boldsymbol{\tau}^n\) is defined as
\[
\boldsymbol{m}(t^{n+1}) + \frac{\Delta t}{2}\left(\boldsymbol{m}(t^{n+1}) \times \Delta_h \boldsymbol{m}(t^{n+1})\right)
= \boldsymbol{m}(t^n) - \frac{\Delta t}{2}\left(\boldsymbol{m}(t^n) \times \Delta_h \boldsymbol{m}(t^n)\right) + \boldsymbol{\tau}^n.
\]
Subtracting the exact equation from the discrete scheme, we get
\[
\boldsymbol{e}^{n+1} + \frac{\Delta t}{2}\left(\boldsymbol{m}^{n+1} \times \Delta_h \boldsymbol{e}^{n+1} + \boldsymbol{e}^{n+1} \times \Delta_h \boldsymbol{m}(t^{n+1})\right)
= \boldsymbol{e}^n - \frac{\Delta t}{2}\left(\boldsymbol{m}^n \times \Delta_h \boldsymbol{e}^n + \boldsymbol{e}^n \times \Delta_h \boldsymbol{m}(t^n)\right) + \boldsymbol{\tau}^n.
\]
Taking the \(L^2\) inner product with \(\boldsymbol{e}^{n+1}\) and using the orthogonality \(\langle \boldsymbol{a} \times \boldsymbol{b}, \boldsymbol{a} \rangle = 0\), we derive
\[
\|\boldsymbol{e}^{n+1}\|_{L^2}^2 \leq \|\boldsymbol{e}^n\|_{L^2}^2 + C\Delta t\left(\|\boldsymbol{\tau}^n\|_{L^2}^2 + \|\boldsymbol{e}^n\|_{L^2}^2\right).
\]
By Gronwall's inequality,
\[
\|\boldsymbol{e}^n\|_{L^2} \leq C e^{CT}\left(\|\boldsymbol{e}^0\|_{L^2} + \sqrt{\Delta t}\sum_{k=0}^{n-1}\|\boldsymbol{\tau}^k\|_{L^2}\right)
\leq C(T)\left(\|\boldsymbol{e}^0\|_{L^2} + \Delta t^2 + h^2\right).
\]
If $\|\boldsymbol{e}^0\|_{L^2} =0$, we have $\|\boldsymbol{e}^n\|_{L^2} = O(\Delta t^2 + h^2)$. By the Lax equivalence theorem, consistency and stability imply convergence. As \(\Delta t, h \to 0\),
\[
\max_{0 \leq n \leq T/\Delta t}\|\boldsymbol{m}^n - \boldsymbol{m}(t^n)\|_{L^2} \to 0,
\]
with second-order accuracy in both time and space:
\[
\|\boldsymbol{m}^n - \boldsymbol{m}(t^n)\|_{L^2} = O(\Delta t^2 + h^2).
\]
\end{remark}


\subsection{Consistency of our proposed scheme}
Considering a scheme below,
\begin{align*}
    \frac{\m_h^n-\m_h^n}{\Delta t}=-\frac{\m_h^{n+1}+\m_h^n}{2}\times \Delta_h \g_h,
\end{align*}
where $\g_h^s=(I-\Delta t \Delta_h)^{-1}\m_h^s$, $s=n,n+1$.

To analyze several specific case, we take
\begin{align*}
    \frac{\m_h^n-\m_h^n}{\Delta t}=-\frac{\m_h^{n+1}+\m_h^n}{2}\times \Delta_h \g_h,
\end{align*}
where $\g_h^n=(I-\Delta t \Delta_h)^{-1}\m_h^n$. The truncation error \(\tau_h^n\) is defined as the residual when substituting the exact solution \(\boldsymbol{m}(t)\) into the discrete scheme:
\[
\tau_h^n = \frac{\boldsymbol{m}(t_{n+1}) - \boldsymbol{m}(t_n)}{\Delta t} + \frac{\boldsymbol{m}(t_{n+1}) + \boldsymbol{m}(t_n)}{2} \times \Delta_h \boldsymbol{g}(t_n),
\]
where \(\boldsymbol{g}(t_n) = (I - \Delta t \Delta_h)^{-1} \boldsymbol{m}(t_n)\), and we assume \(\Delta_h\) is a consistent approximation of the continuous Laplacian \(\Delta\), i.e., \(\Delta_h \boldsymbol{m} = \Delta \boldsymbol{m} + \mathcal{O}(h^2)\) for smooth \(\boldsymbol{m}\).
We expand the exact solution \(\boldsymbol{m}(t_{n+1})\) around \(t_n\) using Taylor's formula:
\[
\boldsymbol{m}(t_{n+1}) = \boldsymbol{m}(t_n) + \Delta t \partial_t \boldsymbol{m}(t_n) + \frac{(\Delta t)^2}{2} \partial_t^2 \boldsymbol{m}(t_n) + \mathcal{O}((\Delta t)^3).
\]
Substituting into the time difference:
\[
\frac{\boldsymbol{m}(t_{n+1}) - \boldsymbol{m}(t_n)}{\Delta t} = \partial_t \boldsymbol{m}(t_n) + \frac{\Delta t}{2} \partial_t^2 \boldsymbol{m}(t_n) + \mathcal{O}((\Delta t)^2).
\]
Similarly,
\[
\frac{\boldsymbol{m}(t_{n+1}) + \boldsymbol{m}(t_n)}{2} = \m(t_n)+\frac{\Delta t}{2} \partial_t \boldsymbol{m}(t_n) + \frac{(\Delta t)^2}{4} \partial_t^2 \boldsymbol{m}(t_n) + \mathcal{O}((\Delta t)^3).
\]
The operator \((I - \Delta t \Delta_h)^{-1}\) can be expanded via Neumann series (for small \(\Delta t\)):
\[
(I - \Delta t \Delta_h)^{-1} = I + \Delta t \Delta_h + (\Delta t)^2 \Delta_h^2 + \mathcal{O}((\Delta t)^3).
\]
Thus,
\[
\boldsymbol{g}(t_n) = \boldsymbol{m}(t_n) + \Delta t \Delta_h \boldsymbol{m}(t_n) + \mathcal{O}((\Delta t)^2+h^2).
\]
Applying \(\Delta_h\) (consistent with \(\Delta\)):
\[
\Delta_h \boldsymbol{g}(t_n) = \Delta_h \boldsymbol{m}(t_n) + \Delta t \Delta_h^2 \boldsymbol{m}(t_n) + \mathcal{O}((\Delta t)^2+h^2) = \Delta \boldsymbol{m}(t_n) + \mathcal{O}(h^2 + \Delta t).
\]
Substitute the expansions into \(\tau_h^n\):
\[
\begin{aligned}
\tau_h^n &= \left( \partial_t \boldsymbol{m}(t_n) + \frac{\Delta t}{2} \partial_t^2 \boldsymbol{m}(t_n) \right) + \left( \m(t_n)+\frac{\Delta t}{2} \partial_t \boldsymbol{m}(t_n) + \mathcal{O}((\Delta t)^2) \right) \times \left( \Delta \boldsymbol{m}(t_n) + \mathcal{O}(h^2 + \Delta t) \right) \\
&\quad = \left[ \partial_t \boldsymbol{m} + \boldsymbol{m} \times \Delta \boldsymbol{m} \right]_{t=t_n}+\mathcal{O}( h^2+\Delta t) \\
&= \mathcal{O}( h^2+\Delta t),
\end{aligned}
\]
where the second term is the exact PDE (assuming the continuous equation is \(\partial_t \boldsymbol{m} = -\boldsymbol{m} \times \Delta \boldsymbol{m}\)).

Considering another specific implicit scheme below,
\begin{align}\label{eq-im}
    \frac{\m_h^{n+1}-\m_h^n}{\Delta t}=-\frac{\m_h^{n+1}+\m_h^n}{2}\times \Delta_h \g_h^{n+1},
\end{align}
where $\g_h^{n+1}=(I-\Delta t \Delta_h)^{-1}\m_h^{n+1}$.
The local truncation error is given by
\[
\tau_h^n = \frac{\boldsymbol{m}(t_{n+1}) - \boldsymbol{m}(t_n)}{\Delta t} + \frac{\boldsymbol{m}(t_{n+1}) + \boldsymbol{m}(t_n)}{2} \times \Delta_h \boldsymbol{g}(t_{n+1}),
\]
Let $\boldsymbol{m}(t)$ be the exact solution of the PDE. Expand $\boldsymbol{m}(t_{n+1})$ around $t_n$ using Taylor series:
\[
\boldsymbol{m}(t_{n+1}) = \boldsymbol{m}(t_n) + \Delta t \frac{\partial \boldsymbol{m}}{\partial t}(t_n) + \frac{(\Delta t)^2}{2} \frac{\partial^2 \boldsymbol{m}}{\partial t^2}(t_n) + \frac{(\Delta t)^3}{6} \frac{\partial^3 \boldsymbol{m}}{\partial t^3}(t_n) + \mathcal{O}((\Delta t)^4),
\]
\[
\boldsymbol{m}(t_{n+1}) + \boldsymbol{m}(t_n) = 2\boldsymbol{m}(t_n) + \Delta t \frac{\partial \boldsymbol{m}}{\partial t}(t_n) + \frac{(\Delta t)^2}{2} \frac{\partial^2 \boldsymbol{m}}{\partial t^2}(t_n) + \mathcal{O}((\Delta t)^3).
\]
For the auxiliary variable $\boldsymbol{g}$, we expand the inverse operator in a Neumann series (valid for small $\Delta t$):
\[
(I - \Delta t \Delta_h)^{-1} = I + \Delta t \Delta_h + (\Delta t)^2 \Delta_h^2 + \mathcal{O}((\Delta t)^3),
\]
thus:
\[
\boldsymbol{g} = \boldsymbol{m} + \Delta t \Delta_h \boldsymbol{m} + (\Delta t)^2 \Delta_h^2 \boldsymbol{m} + \mathcal{O}((\Delta t)^3),
\]
\[
\Delta \boldsymbol{g} = \Delta \boldsymbol{m} + \Delta t \Delta_h^2 \boldsymbol{m} + \mathcal{O}((\Delta t)^2).
\]
Substitute the Taylor expansions into the definition of $\tau^n$:

First, rewrite the time difference term:
\[
\frac{\boldsymbol{m}(t_{n+1}) - \boldsymbol{m}(t_n)}{\Delta t} = \frac{\partial \boldsymbol{m}}{\partial t}(t_n) + \frac{\Delta t}{2} \frac{\partial^2 \boldsymbol{m}}{\partial t^2}(t_n) + \frac{(\Delta t)^2}{6} \frac{\partial^3 \boldsymbol{m}}{\partial t^3}(t_n) + \mathcal{O}((\Delta t)^3).
\]

Second, rewrite the average term:
\[
\frac{\boldsymbol{m}(t_{n+1}) + \boldsymbol{m}(t_n)}{2} = \boldsymbol{m}(t_n) + \frac{\Delta t}{2} \frac{\partial \boldsymbol{m}}{\partial t}(t_n) + \frac{(\Delta t)^2}{4} \frac{\partial^2 \boldsymbol{m}}{\partial t^2}(t_n) + \mathcal{O}((\Delta t)^3).
\]

Third, expand $\Delta_h \boldsymbol{g}(t_{n+1})$ (assuming $\Delta_h$ is consistent with $\Delta$):
\[
\Delta_h \boldsymbol{g}(t_{n+1}) = \Delta_h \boldsymbol{m}(t_{n+1}) + \Delta t \Delta_h^2 \boldsymbol{m}(t_{n+1}) + \mathcal{O}((\Delta t)^2+h^2) = \Delta_h \boldsymbol{m}(t_n)  + \mathcal{O}(\Delta t+h^2).
\]

Substitute these into $\tau^n$:
\begin{align*}
    \tau_h^n &= \left( \frac{\partial \boldsymbol{m}}{\partial t}(t_n) + \frac{\Delta t}{2} \frac{\partial^2 \boldsymbol{m}}{\partial t^2}(t_n) + \mathcal{O}((\Delta t)^2) \right) + \left( \boldsymbol{m}(t_n) + \frac{\Delta t}{2} \frac{\partial \boldsymbol{m}}{\partial t}(t_n) + \mathcal{O}((\Delta t)^2) \right) \times \left( \Delta_h \boldsymbol{m}(t_n) +  \mathcal{O}(\Delta t+h^2) \right)\\
    &=(\partial_t \m+\m\times \Delta_h\m)|_{t=t_n}+O(\Delta t+h^2).
\end{align*}

Now we consider our proposed method, and note that
\begin{align*}
    \frac{\m_h^n-\m_h^n}{\Delta t}=-\frac{\m_h^{n+1}+\m_h^n}{2}\times \Delta_h \g_h^s,
\end{align*}
where $\g_h^{s}=(I-\Delta t \Delta_h)^{-1}\m_h^{s}$, $s=n,n+1$. Such a method with a Gauss-Seidel type iteration to update the $g_h$. From previous local truncation analysis for explicit and implicit schemes, we can directly derive that
\begin{align*}
    \tau_h^n=O(\Delta t+h^2).
\end{align*}



\subsection{Stability of our proposed scheme}
For the step 1, we denote $A$ for its coefficient matrix, and note that the three eigenvalue for $A_i$ and $A_i^{-1}$ are all $1$, so the step 1 is stable.

For step 2, we denote $A$ for its coefficient matrix without any generality. To be specific,
\begin{align}
    A=\begin{pmatrix}
        1 & c_n&-b_{n}\\
        -c_n &1 &a_{n+1}\\
        0 &0&1
    \end{pmatrix}
\end{align}
Note that $\det(A-\lambda I)=(1-\lambda)[(1-\lambda)^2+c_n^2]$, we have the eigenvalues of $A$ below,
\begin{align*}
    \lambda_1=1,\quad \lambda_{2,3}=1\pm ic_n.
\end{align*}
Then $|\lambda_1|=1$, $|\lambda_2|=|\lambda_3|=\sqrt{1+c_n^2}$. Thus, the three eigenvalues of $A_i^{-1}$ are 
\begin{align*}
    \gamma_1=1/\lambda_1,\quad \gamma_2=1/\lambda_2,\quad \gamma_3=1/\lambda_3.
\end{align*}
The absolute values of three eigenvalues of $A_i^{-1}$ are 
\begin{align*}
    |\gamma_1|=1,\quad |\gamma_2|=|\gamma_3|=\frac{1}{\sqrt{1+c_n^2}}\le 1.
\end{align*}
Here $c_n$ is small compared to $1$.

For step 3, we denote $A$ for its coefficient matrix without any generality. To be specific,
\begin{align}
    A=\begin{pmatrix}
        1 & c_n&-b_{n+1}\\
        -c_n &1 &a_{n+1}\\
        b_{n+1} &-a_{n+1} &1
    \end{pmatrix}
\end{align}
Note that
\begin{align*}
    \det(A-\lambda I)=(1-\lambda)^3+(1-\lambda)(a_{n+1}^2+b_{n+1}^2+c_n^2)=(1-\lambda)[(1-\lambda)^2+a_{n+1}^2+b_{n+1}^2+c_n^2].
\end{align*}
Thus the eigenvalues of $A$ are
\begin{align*}
    \lambda_1=1,\quad \lambda_2=1+i\sqrt{a_{n+1}^2+b_{n+1}^2+c_n^2},\quad \lambda_3=1-i\sqrt{a_{n+1}^2+b_{n+1}^2+c_n^2}.
\end{align*}
Then $|\lambda_1|=1$, $|\lambda_2|=|\lambda_3|=\sqrt{1+a_{n+1}^2+b_{n+1}^2+c_n^2}$. Thus, the three eigenvalues of $A_i^{-1}$ are 
\begin{align*}
    \gamma_1=1/\lambda_1,\quad \gamma_2=1/\lambda_2,\quad \gamma_3=1/\lambda_3.
\end{align*}
The absolute values of three eigenvalues of $A_i^{-1}$ are 
\begin{align*}
    |\gamma_1|=1,\quad |\gamma_2|=|\gamma_3|=\frac{1}{\sqrt{1+a_{n+1}^2+b_{n+1}^2+c_n^2}}\le 1.
\end{align*}
Here $a_{n+1},b_{n+1},c_n$ is small compared to $1$.

\begin{remark}
We first define explicit and implicit methods using the general form of time integration for a differential equation:
\[
\frac{du}{dt} = f(t, u), \quad u(t_0) = u_0
\]
where $u(t)$ is the solution variable, $t$ is time, and $f(\cdot)$ is the right-hand side (RHS) function.

Explicit methods compute the next time step $u_{n+1}$ using only \textit{known} values from previous steps ($u_n, u_{n-1}, \dots$):
\[
u_{n+1} = u_n + \Delta t \cdot F_{\text{explicit}}(t_n, u_n)
\]
The most common example is the \textbf{Forward Euler Method}:
\[
u_{n+1} = u_n + \Delta t \cdot f(t_n, u_n)
\]
where $\Delta t = t_{n+1} - t_n$ is the time step size.

Implicit methods compute $u_{n+1}$ using \textit{unknown} values at the next time step ($u_{n+1}$):
\[
u_{n+1} = u_n + \Delta t \cdot F_{\text{implicit}}(t_{n+1}, u_{n+1})
\]
The classic example is the \textbf{Backward Euler Method}:
\[
u_{n+1} = u_n + \Delta t \cdot f(t_{n+1}, u_{n+1})
\]
This requires solving an algebraic equation (linear/nonlinear) for $u_{n+1}$ at each step.

Stability is the most critical advantage of implicit methods. We illustrate this with the linear test problem:
\[
\frac{du}{dt} = \lambda u, \quad \lambda \in \mathbb{C}, \text{Re}(\lambda) < 0
\]

Applying Forward Euler gives:
\[
u_{n+1} = (1 + \lambda \Delta t) u_n
\]
For stability ($|u_{n+1}| \leq |u_n|$), we require:
\[
|1 + \lambda \Delta t| \leq 1
\]
This imposes a \textit{strict upper bound} on $\Delta t$ (e.g., for $\lambda = -10$, $\Delta t < 0.2$). For stiff problems (large $|\lambda|$), $\Delta t$ must be extremely small.

Applying Backward Euler gives:
\[
u_{n+1} = \frac{1}{1 - \lambda \Delta t} u_n
\]
For stability:
\[
\left| \frac{1}{1 - \lambda \Delta t} \right| \leq 1
\]
This holds for \textit{any positive} $\Delta t$ when $\text{Re}(\lambda) < 0$ (A-stable). Implicit methods have \textit{unconditional stability} for stiff problems, while explicit methods have \textit{conditional stability}.
\end{remark}

For our step 3, the matrix 
\begin{align*}
    A=\begin{pmatrix}
        1 & c &-b\\
        -c & 1 & a\\
        b & -a & 1
    \end{pmatrix}
\end{align*}
Here are more comparison details for stability:
\begin{itemize}
    \item For Crank-Nicolson's scheme, $a, b, c$ directly depends on $\m^{n+1}$, so the implicit way gives great stability (A-stability). However, it has to be a nonlinear system.
    \item If $a,b,c$ directly depends on $\m^n$, we get an explicit scheme with CFL-type conditional stability, even though the lower order term is implicit.
    \item If $a,b,c$ indirectly depends on $\m^n$, say implicit treatment by solving a heat diffusion equation, we get a slightly better CFL-condition.
    \item Our method design a Gauss-Seidel type iteration to update the $m^{n+1}$, thus $a, b,c$ can be formed as a mixed way to combine $\m^n$ and $\m^{n+1}$. Such a method gives unconditionally stable results.
\end{itemize}

\subsection{Convergence of our proposed scheme}
By the Lax equivalence theorem, consistency and stability imply convergence. As \(\Delta t, h \to 0\),
\[
\max_{0 \leq n \leq T/\Delta t}\|\boldsymbol{m}^n - \boldsymbol{m}(t^n)\|_{L^2} \to 0,
\]
with first-order accuracy in time and second-order accuracy in space:
\[
\|\boldsymbol{m}^n - \boldsymbol{m}(t^n)\|_{L^2} = O(\Delta t + h^2).
\]

Considering the implicit scheme \cref{eq-im}, we denote $e_h^n=\m_h^n-\m_h(t_n)$, and have

\begin{align*}
    \frac{e_h^{n+1}-\e_h^n}{\Delta t}=-\frac{\e_h^{n+1}+\e_h^n}{2}\times \Delta_h \g_h^{n+1}-\frac{\m(t_{n+1})+\m(t_n)}{2}\times \Delta_h (I-\Delta t\Delta_h)^{-1}\e_h^{n+1}+O(\Delta t+h^2)
\end{align*}
Taking the inner product with $\e_h^{n+1}+\e_h^n$, we have
\begin{align*}
    \|\e_h^{n+1}\|^2-\|\e_h^n\|^2=-\Delta t \langle \frac{\m(t_{n+1})+\m(t_n)}{2}\times \Delta_h (I-\Delta t\Delta_h)^{-1}\e_h^{n+1},\e_h^{n+1}+\e_h^n \rangle +O(\Delta t+h^2).
\end{align*}
Indeed, we have
\[
\|e_h^n\|_{L^2} \leq C (\Delta t + h^2),\quad \forall n\in \mathbb{Z}
\]
which will be proved in the future work.

\begin{lemma}[Boundedness of Elliptic Operator]
For any $\Delta t > 0$ and mesh size $h > 0$, the operator $(I - \Delta t \Delta_h)^{-1}$ is uniformly bounded in $L^2(\Omega)$:
\[
\|(I - \Delta t \Delta_h)^{-1} u\|_{L^2} \leq C \|u\|_{L^2},
\]
where $C$ is a constant independent of $\Delta t$ and $h$. Moreover,
\[
\|\Delta_h (I - \Delta t \Delta_h)^{-1} u\|_{L^2} \leq C \|u\|_{L^2}.
\]
\end{lemma}

\begin{proof}
Since $-\Delta_h$ is positive definite with spectral bound $\sigma(-\Delta_h) \leq C h^{-2}$, we have
\[
\|(I - \Delta t \Delta_h)^{-1}\| = \sup_{\|u\|=1} \frac{1}{\|(I - \Delta t \Delta_h)u\|} \leq \frac{1}{1 + \Delta t \cdot \inf\sigma(-\Delta_h)} \leq C.
\]
The boundedness of $\Delta_h (I - \Delta t \Delta_h)^{-1}$ follows from $\|\Delta_h (I - \Delta t \Delta_h)^{-1}\| = \|-(I - \Delta t \Delta_h)^{-1}(-\Delta_h)\| \leq C$.
\end{proof}

% Using the Cauchy-Schwarz inequality and $L^\infty$-boundedness of $m(t)$ ($\|m(t)\|_{L^\infty} \leq M$):
% \[
% \left| \left\langle \frac{\m(t_{n+1}) + \m(t_n)}{2} \Delta_h (I - \Delta t \Delta_h)^{-1} \e_h^{n+1}, \e_h^{n+1}+\e_h^n \right\rangle \right| \leq M \cdot \|\Delta_h (I - \Delta t \Delta_h)^{-1}\| \cdot \left(\frac32\|\e_h^{n+1}\|_{L^2}^2 +\frac12\|\e_h^{n}\|_{L^2}^2 \right)\leq C (\|\e_h^{n+1}\|_{L^2}^2+\|\e_h^{n}\|_{L^2}^2).
% \]

% Rearranging the error equation, we have
% \[
% \|e_h^{n+1}\|_{L^2}^2 \leq (1-C\Delta t)^{-1}(1 + C\Delta t) \|e_h^n\|_{L^2}^2 + C \Delta t (\Delta t + h^2).
% \]
% Applying the discrete Gronwall's inequality:
% \[
% \|e_h^N\|_{L^2}^2 \leq e^{CT} \|e_h^0\|_{L^2}^2 + C T (\Delta t^2 + h^4).
% \]
% Assuming the initial error satisfies $\|e_h^0\|_{L^2} = \mathcal{O}(h^2)$ (second-order spatial approximation), we get:



\section{Numerical experiments}
\label{sec:experiments}

In this section, we proceed the accuracy, stability and norm preserving test. We choose different initial conditions to have the robustness of such a method.

\subsection{Accuracy and efficiency tests}

To simplify the accuracy verification, we set the exact solution given for the governing model \cref{eq-5}. Analytical exact solutions are derived for both one-dimensional (1D) and three-dimensional (3D) scenarios to serve as benchmarks for error quantification.

For the 1D case, the exact magnetization solution \(\m_e\) is:
\begin{equation*}
\m_e=\left(\cos(\cos(\pi x))\sin t, \sin(\cos(\pi x))\sin t, \cos t\right)^T,
\end{equation*}
while the corresponding 3D exact solution is:
\begin{equation*}
\m_e=\left(\cos(XYZ)\sin t, \sin(XYZ)\sin t, \cos t\right)^T,
\end{equation*}
where $X=x^2(1-x)^2$, $Y=y^2(1-y)^2$ and $Z=z^2(1-z)^2$.

These exact solutions satisfy the governing equation \cref{eq-5} when the forcing term is defined as \(\f_e=\partial_t \m_e+\m_e \times \Delta \m_e+\alpha\m_e\times (\m_e \times \Delta \m_e)\). They also comply with the homogeneous Neumann boundary condition, ensuring consistency with simulation constraints.

To isolate the temporal approximation error from spatial discretization effects, the spatial resolution in the 1D test is fixed at \(h=5\times 10^{-4}\)—a sufficiently fine grid that makes spatial error negligible compared to temporal error. The Gilbert damping parameter is set to \(\alpha=0.01\), and simulations run until final time \(T=0.1\). Under this configuration, the measured error primarily reflects temporal discretization inaccuracy.

The 3D temporal accuracy test faces inherent constraints from spatial resolution, as excessively fine grids incur prohibitive computational cost. To balance spatial and temporal error contributions, we adopt a coordinated refinement strategy for spatial mesh sizes (\(h_x, h_y, h_z\)) and temporal step-size (\(\Delta t\)) tailored to the proposed method's order: \(\Delta t=h_x^2=h_y^2=h_z^2=h^2=T/N_0\). Here, \(N_0\) is a refinement level parameter, with specific values given in subsequent results. Consistent with the 1D test, \(\alpha=0.01\), and the final time \(T\) is specified later. The first order temporal accuracy is verified from \cref{tab-a-v-3} and \cref{tab-a-10-time-3D-Q-2} for 1D and 3D tests, respectively. 

\begin{table}[htbp]
	\centering
	\caption{The temporal accuracy in 1D test for proposed method with damping $\alpha=0.01$when $h = 5D-4$, $T=1d-1$.}\label{tab-a-v-3}
	\begin{tabular}{|c|c|c|c|}
		\hline
		$k$ & $\|\m_h-\m_e\|_\infty$ & $\|\m_h-\m_e\|_2$ &$\|\m_h-\m_e\|_{H^1}$\\
		\hline
		$T/80$ &0.001304094971804&8.500523347099678e-04&0.006116653503286 \\
		$T/120$ &8.684032607750442e-04&5.745020596842719e-04&0.004123837837786 \\
		$T/160$ &6.505721097687933e-04&4.340035210414707e-04&0.003112167088311 \\
		$T/240$ &4.330118558566187e-04&2.915444972028109e-04&0.002089104867675 \\
		$T/320$ &3.244330910497223e-04&2.195304920518377e-04&0.001572755516106 \\
		\hline 
		order &1.003609279663207&0.976857797102320&0.979916584127108\\
		\hline
	\end{tabular}
\end{table}

\begin{table}[htbp]
	\centering
	\caption{The spatial accuracy in 1D test for proposed method with damping $\alpha=0.01$ when $k= 1D-6$, $T=1d-1$.}\label{tab-a-10-space-Q-1}
	\begin{tabular}{|c|c|c|c|}
		\hline
		$h$ & $\|\m_h-\m_e\|_\infty$ & $\|\m_h-\m_e\|_2$ &$\|\m_h-\m_e\|_{H^1}$ \\
		\hline
		1/16 &4.225596750053739e-04&2.896508432807531e-04&0.002209985483017 \\
		1/24 &1.885253776899853e-04&1.286680596306939e-04&9.768130592826686e-04\\
		1/32 &1.062644247209338e-04&7.252233752489178e-05&5.480246619562391e-04\\
		1/48 &4.739270964135289e-05
		&3.248282203669158e-05&2.426999432605687e-04\\
		1/64 &2.676411577153676e-05&1.848377745692578e-05&1.360493878062051e-04\\
		\hline 
		order &1.990738385102109&1.985405237927322&2.010529053514131 \\
		\hline
	\end{tabular}
\end{table}

Following temporal accuracy evaluation, spatial accuracy tests were conducted to quantify the spatial discretization performance of the proposed method. To prevent temporal errors from interfering with spatial accuracy assessment, the temporal step size was fixed at a sufficiently small \(k=10^{-6}\) for 1D test shown in \cref{tab-a-10-space-Q-1}—making temporal errors negligible compared to spatial errors. We take $k=h^2$, such that the second order spatial accuracy is observed, which is consistent with 1D test.

\begin{table}[htbp]
	\centering
	\caption{The temporal accuracy and spatial accuracy for the proposed method with damping $\alpha=0.01$, $T=0.1$ with $k=h^2$ in 3D. }\label{tab-a-10-time-3D-Q-2}
	\begin{tabular}{|c|c|c|c|c|}
		\hline
		$k$&$h$ & $\|\m_h-\m_e\|_\infty$ & $\|\m_h-\m_e\|_2$ &$\|\m_h-\m_e\|_{H^1}$ \\
	\hline
	T/10&1/10&5.006365255465495e-04&2.886424573026357e-04&3.499752082858884e-04 \iffalse 3.435768 (s) \fi\\
	T/40 &1/20&1.264524159770852e-04&7.243063478581653e-05&1.290889967619724e-04 \iffalse 54.921250 (s) \fi\\
	T/57 &1/24&8.912305801656029e-05&5.101958217642231e-05&1.062422830734351e-04\iffalse 127.851179 (s) \fi\\
	T/78&1/28&6.546837939636063e-05&3.751573538129243e-05&9.169167058822856e-05 \iffalse 285.484803 (s) \fi\\
	T/102 &1/32&5.037225571857817e-05&2.895218634093047e-05&8.242737703221056e-05 \iffalse 581.983777 (s) \fi\\
	% T/129 &1/36& &&\\
	\hline 
	order &&0.989524179094233&0.991610312809092&0.635976416720683\\
		\hline
	&	order &1.973909640037754&1.978077601727147&1.268736008455239\\
		\hline
	\end{tabular}
\end{table}

\subsection{Norm preserving tests}
We choose the initial condition with below,
\begin{align*}
    \m_0=\left(\cos(\cos(\pi x))\sin (0.01), \sin(\cos(\pi x))\sin (0.01), \cos (0.01)\right)^T,
\end{align*}
and \begin{align*}
\m_0=\left(\cos(XYZ)\sin (0.01), \sin(XYZ)\sin (0.01), \cos (0.01)\right)^T,
\end{align*}
in 1D and 3D, respectively. The results are presented in \cref{tab-8} and \cref{tab-9}.

\begin{table}[htbp]
	\centering
	\caption{The proposed method with damping $\alpha=0.01$ when $h = 5D-4$, $T=1d-1$ for the norm preserving test in 1D.}\label{tab-8}
	\begin{tabular}{|c|c|}
		\hline
		$k$ & $\|\|\m_h\|_2-1\|_\infty$ \\
		\hline
	    2.0D-2 & 1.110223024625157e-15 \\
		1.0D-2 & 2.331468351712829e-15 \\
		5.0D-3 & 2.886579864025407e-15\\
		2.5D-3 & 3.996802888650564e-15 \\
		1.25D-3 & 5.995204332975845e-15 \\
		6.25D-4 & 8.881784197001252e-15 \\
		3.125D-4 & 1.165734175856414e-14\\
		\hline
	\end{tabular}
\end{table}


\begin{table}[htbp]
	\centering
	\caption{The norm preserving test for the proposed method with damping $\alpha=0.01$, $T=0.1$ with $k=h^2$ in 3D. }\label{tab-9}
	\begin{tabular}{|c|c|c|}
		\hline
		$k$&$h$ & $\|\|\m_h\|_2-1\|_\infty$  \\
	\hline
	T/10&1/10& 5.748734821509061e-13\\
	T/40 &1/20&4.660716257376407e-13 \\
	T/57 &1/24&4.194422587033841e-13 \\
	T/78&1/28& 3.410605131648481e-13\\
	% T/102 &1/32&\\
	% T/129 &1/36&  \\
	\hline 
	\end{tabular}
\end{table}

\subsection{Initial conditions and stability tests}

In this section, we choose different initial conditions specified in each test.

In 1D, we specify the initial condition as
\begin{align*}
    \m_0&=\left(\cos(x^2(1-x)^2)\sin (0.01), \sin(x^2(1-x)^2)\sin (0.01), \cos (0.01)\right)^T\\
    \m_0&=\left(\cos(\cos(\pi x))\sin (0.01), \sin(\cos(\pi x))\sin (0.01), \cos (0.01)\right)^T.
\end{align*}

The results for 1D case using the proposed method to show the numerical profiles are presented in \Cref{fig:1} with the parameters $\alpha=0.01$, $N_x=2000$ and $N_t=5$. 

In 3D, we choose the initial conditions below,
\begin{align*}
\m_0&=\left(\cos(x^2(1-x)^2)\sin (0.01), \sin(x^2(1-x)^2)\sin (0.01), \cos (0.01)\right)^T\\
    \m_0&=\left(\cos(\cos(\pi x))\sin (0.01), \sin(\cos(\pi x))\sin (0.01), \cos (0.01)\right)^T\\
    % \m_0&=\left(\cos(\cos(\pi x))\sin (x+t), \sin(\cos(\pi x))\sin (x+t), \cos (x+t)\right)^T\\
    \m_0&=\left(\cos(\cos(\cos(\pi x)))\sin (\pi x+0), \sin(\cos(\cos(\pi x)))\sin (\pi x+0), \cos (\pi x+0)\right)^T\\
    %  \m_0&=\left(\cos(X)\sin (0.01), \sin(X)\sin (0.01), \cos (0.01)\right)^T.
\end{align*}
The results for those initial conditions are presented in \Cref{fig:2-2}, \Cref{fig:2} and \Cref{fig:2-1} which verify the consistency.

% To get the robustness, we test other initial conditions below,
% \begin{align*}
%     % \m_0&=\left(\cos(\tan(\pi x))\sin (0.01), \sin(\tan(\pi x))\sin (0.01), \cos (0.01)\right)^T\\
%     %  \m_0&=\left(\cos(2(x+y+z))\sin (0.01), \sin(2(x+y+z))\sin (0.01), \cos (0.01)\right)^T\\
%     %    \m_0&=\left(\cos(\cos(\pi x)\cos(\pi y))\sin (0.01), \sin(\cos(\pi x)\cos(\pi y))\sin (0.01), \cos (0.01)\right)^T\\
%         \m_0&=\left(\cos(XYZ)\sin (0.01), \sin(XYZ)\sin (0.01), \cos (0.01)\right)^T\\
%          \m_0&=\left(\cos(\cos(\pi x)\cos(\pi y)\cos(\pi z))\sin (0.01), \sin(\cos(\pi x)\cos(\pi y)\cos(\pi z))\sin (0.01), \cos (0.01)\right)^T.
% \end{align*}
% The results for other initial conditions are presented in \Cref{fig:3} and \Cref{fig:4} which verify the stability and robustness.

\begin{figure}[htbp]
    \centering
    \subfloat[$m_1$]{\includegraphics[width=0.5\linewidth]{v1/m1_1D_v1_v1.png}}
    \subfloat[$m_1$]{\includegraphics[width=0.5\linewidth]{v1/m1_1D_v2.png}}
    \hspace{0.1in}
    \subfloat[$m_2$]{\includegraphics[width=0.5\linewidth]{v1/m2_1D_v1_v1.png}}
    \subfloat[$m_2$]{\includegraphics[width=0.5\linewidth]{v1/m2_1D_v2.png}}
    \hspace{0.1in}
    \subfloat[$m_3$]{\includegraphics[width=0.5\linewidth]{v1/m3_1D_v1_v1.png}}
    \subfloat[$m_3$]{\includegraphics[width=0.5\linewidth]{v1/m3_1D_v2.png}}
    \caption{The solution profile using proposed method in 1D given the initial condition $m_0$ with $T0$ specified without source term, $\alpha=0.01$ and $T=0.1$, $N_x=2000$, $N_t=5$.}
    \label{fig:1}
\end{figure}

\begin{figure}[htbp]
    \centering
    \subfloat[arrow profile]{\includegraphics[width=0.5\linewidth]{v1/initial_8_3D_arrow_v2.png}}
    \subfloat[angle profile]{\includegraphics[width=0.53\linewidth]{v1/initial_8_3D_color_v2.png}}
    \hspace{0.1in}
    \subfloat[arrow profile]{\includegraphics[width=0.5\linewidth]{v1/proposed_8_3D_arrow_v2.png}}
    \subfloat[angle profile]{\includegraphics[width=0.53\linewidth]{v1/proposed_8_3D_color_v2.png}}
      \caption{The solution profile using the proposed methods in 3D given the initial condition $m_0$ with initial condition specified without source term, $\alpha=0$ and $T=0.1$, $N_x=N_y=N_z=20$, $N_t=40$. Top row with initial condition; Bottom row with proposed method. Initial condition given: $\m_0=[\cos(x^2(1-x)^2)\sin(0.01),\sin(x^2(1-x)^2)\sin(0.01),\cos(0.01)]$.}
    \label{fig:2-2}
\end{figure}

\begin{figure}[htbp]
    \centering
    \subfloat[arrow profile]{\includegraphics[width=0.5\linewidth]{v1/initial_1_3D_arrow_v1.png}}
    \subfloat[angle profile]{\includegraphics[width=0.53\linewidth]{v1/initial_1_3D_color_v1.png}}
    \hspace{0.1in}
    \subfloat[arrow profile]{\includegraphics[width=0.5\linewidth]{v1/proposed_3D_arrow_initial_1_v1.png}}
    \subfloat[angle profile]{\includegraphics[width=0.53\linewidth]{v1/proposed_3D_color_initial_1_v1.png}}
      \caption{The solution profile using the proposed methods in 3D given the initial condition $m_0$ with initial condition specified without source term, $\alpha=0$ and $T=0.1$, $N_x=N_y=N_z=20$, $N_t=40$. Top row with initial condition; Bottom row with proposed method. Initial condition given: $\m_0=[\cos(\cos(\pi x))\sin(0.01),\sin(\cos(\pi x))\sin(0.01),\cos(0.01)]$.}
    \label{fig:2}
\end{figure}

\begin{figure}[htbp]
    \centering
    \subfloat[arrow profile]{\includegraphics[width=0.5\linewidth]{v1/initial_1_3D_arrow_v2.png}}
    \subfloat[angle profile]{\includegraphics[width=0.53\linewidth]{v1/initial_1_3D_color_v2.png}}
    \hspace{0.1in}
    \subfloat[arrow profile]{\includegraphics[width=0.5\linewidth]{v1/proposed_3D_arrow_initial_1_v2.png}}
    \subfloat[angle profile]{\includegraphics[width=0.53\linewidth]{v1/proposed_3D_color_initial_1_v2.png}}
      \caption{The solution profile using the proposed methods in 3D given the initial condition $m_0$ with initial condition specified without source term, $\alpha=0$ and $T=0.1$, $N_x=N_y=N_z=20$, $N_t=40$. Top row with initial condition; Bottom row with proposed method. Initial condition given: $\m_0=[\cos(\cos(\cos(\pi x)))\sin(\pi x+T0),\sin(\cos(\cos(\pi x)))\sin(\pi x+T0),\cos(\pi x+T0)]$.}
    \label{fig:2-1}
\end{figure}

% \begin{figure}[htbp]
%     \centering
%     \subfloat[arrow profile]{\includegraphics[width=0.5\linewidth]{v1/initial_3D_arrow_v1.png}}
%     \subfloat[angle profile]{\includegraphics[width=0.53\linewidth]{v1/initial_3D_color_v1.png}}
%     \hspace{0.1in}
%      \subfloat[arrow profile]{\includegraphics[width=0.5\linewidth]{v1/proposed_arrow_3D_v1.png}}
%     \subfloat[angle profile]{\includegraphics[width=0.53\linewidth]{v1/proposed_color_3D_v1.png}}
%       \caption{The magnetization profile using the proposed method given initial conditions in 3D as $\m_0=[\cos(XYZ)\sin(0.01),\sin(XYZ)\sin(0.01),\cos(0.01)]$ specified without source term, $\alpha=0.01$ and $T=0.1$, $N_x=N_y=N_z=20$, $N_t=40$.}
%     \label{fig:3}
% \end{figure}

% \begin{figure}[htbp]
%     \centering
%     % \subfloat[arrow profile]{\includegraphics[width=0.5\linewidth]{v1/initial_3D_arrow_v1.png}}
%     % \subfloat[angle profile]{\includegraphics[width=0.53\linewidth]{v1/initial_3D_color_v1.png}}
%     % \hspace{0.1in}
%     \subfloat[arrow profile]{\includegraphics[width=0.5\linewidth]{v1/initial_3D_arrow_v2.png}}
%     \subfloat[angle profile]{\includegraphics[width=0.53\linewidth]{v1/initial_3D_color_v2.png}}
%     \hspace{0.1in}
%     \subfloat[arrow profile]{\includegraphics[width=0.5\linewidth]{v1/proposed_arrow_3D_v2.png}}
%     \subfloat[angle profile]{\includegraphics[width=0.53\linewidth]{v1/proposed_color_3D_v2.png}}
%       \caption{The magnetization profile using the proposed method given initial conditions in 3D as $\m_0=[\cos(\cos(\pi x)\cos(\pi y)\cos(\pi z))\sin(0.01),\sin(\cos(\pi x)\cos(\pi y)\cos(\pi z))\sin(0.01),\cos(0.01)]$ specified without source term, $\alpha=0.01$ and $T=0.1$, $N_x=N_y=N_z=20$, $N_t=40$.}
%     \label{fig:4}
% \end{figure}


% \begin{figure}[htbp]
%     \centering
%     \subfloat[arrow profile]{\includegraphics[width=0.5\linewidth]{v1/proposed_arrow_3D_v1.png}}
%     \subfloat[angle profile]{\includegraphics[width=0.53\linewidth]{v1/proposed_color_3D_v1.png}}
%     \hspace{0.1in}
%     \subfloat[arrow profile]{\includegraphics[width=0.5\linewidth]{v1/proposed_arrow_3D_v2.png}}
%     \subfloat[angle profile]{\includegraphics[width=0.53\linewidth]{v1/proposed_color_3D_v2.png}}
%       \caption{The solution profile using proposed method in 3D given the initial condition $m_0$ with initial condition specified without source term, $\alpha=0.01$ and $T=0.1$, $N_x=N_y=N_z=20$, $N_t=40$. Top row: $\m_0=[\cos(XYZ)\sin(0.01),\sin(XYZ)\sin(0.01),\cos(0.01)]$; Bottom row: $\m_0=[\cos(\cos(\pi x)\cos(\pi y)\cos(\pi z))\sin(0.01),\sin(\cos(\pi x)\cos(\pi y)\cos(\pi z))\sin(0.01),\cos(0.01)]$}
%     \label{fig:4}
% \end{figure}



% \subsection{Stability tests}

% \subsection{Magnetic switching}

\section{Conclusions and discussions}
\label{sec:conclusions}

In this paper, we propose a structure preserving method with the first order accuracy in time and the second order accuracy in space. Such a method use a Crank-Nicolson's implicit method with the Gauss-Seidel iteration. Such a method preserves the norm constraint. Such a method is constructed only based on the equation, not using a projection step. Implicit Crank-Nicolson's method involves a nonlinear system, however, our proposed method is linear and stable. In the future work, we will analyze the stability of such a method and apply for micromagnetics simulations. We will modify such a method to be a finite element methods for the space discretization. The proposed method is stable, length preserving, accurate. Meanwhile, the convergence analysis and stability analysis for the proposed method along with the normalizing step can be proved. The implicit Crank-Nicolson's method has great theoretical value, such as stability, norm preserving, however, our approach is linear and expected to have some of the good properties of the implicit Crank-Nicolson's method.

\section*{Acknowledgments}
This work is supported in part by the Jiangsu Science and Technology Programme-Fundamental Research Plan Fund (BK20250468), Research and
Development Fund of Xi'an Jiaotong Liverpool University (RDF-24-01-015).


\vspace{1cm}

\bibliographystyle{elsarticle-num-names}
\bibliography{references.bib}
%\begin{thebibliography}{29}
%\bibitem[{Latorre and Riera(2009)}]{Latorre2009}
%{J.~I. Latorre}, {A.~Riera}, \textit{A
%  short review on entanglement in quantum spin systems}, {J.
%  Phys. A: Math. Theor.} {42} ({2009})
%  {504002}.
%
%\bibitem[{Eisert et~al.(2010)Eisert, Cramer, and Plenio}]{Eisert2010}
%{J.~Eisert}, {M.~Cramer},
%  {M.~B. Plenio}, \textit{\textit{Colloquium}: Area
%  laws for the entanglement entropy}, {Rev. Mod. Phys.}
%  {82} ({2010}) {277}.
%  %, \urlprefix\url{http://link.aps.org/doi/10.1103/RevModPhys.82.277}.
%
%\bibitem[{Laflorencie(2016)}]{Laflorencie2016}
%{N.~Laflorencie}, \textit{Quantum entanglement in
%  condensed matter systems}, {Phys. Rep.}
%  {646} ({2016}) {1}.
%  %, \urlprefix\url{http://www.sciencedirect.com/science/article/pii/S0370157316301582}.
%
%\bibitem[{Gioev and Klich(2006)}]{Klich2006}
%{D.~Gioev}, {I.~Klich},
%  \textit{Entanglement Entropy of Fermions in Any Dimension and the
%  {Widom} Conjecture}, {Phys. Rev. Lett.} {96}
%  ({2006}) {100503}.
%  %, \urlprefix\url{http://link.aps.org/doi/10.1103/PhysRevLett.96.100503}.
%
%\bibitem[{Li et~al.(2006)Li, Ding, Yu, Roscilde, and Haas}]{Li2006}
%{W.~Li}, {L.~Ding}, {R.~Yu},
%  {T.~Roscilde}, {S.~Haas},
%  \textit{Scaling behavior of entanglement in two- and
%  three-dimensional free-fermion systems}, {Phys. Rev. B}
%  {74} ({2006}) {073103}.
%  % , \urlprefix\url{http://link.aps.org/doi/10.1103/PhysRevB.74.073103}.
%
%\bibitem[{Swingle(2010)}]{Swingle2010}
%{B.~Swingle}, \textit{Entanglement Entropy and the
%  {Fermi} Surface}, {Phys. Rev. Lett.} {105}
%  ({2010}) {050502}.
%  %,  \urlprefix\url{http://link.aps.org/doi/10.1103/PhysRevLett.105.050502}.
%
%\bibitem[{Swingle(2012)}]{Swingle2012}
%{B.~Swingle}, \textit{R\'enyi entropy, mutual
%  information, and fluctuation properties of {Fermi} liquids},
%  {Phys. Rev. B} {86} ({2012})
%  {045109}.
%  %, \urlprefix\url{http://link.aps.org/doi/10.1103/PhysRevB.86.045109}.
%
%\bibitem[{Klich and Levitov(2009)}]{Klich2009}
%{I.~Klich}, {L.~Levitov},
%  \textit{Quantum Noise as an Entanglement Meter},
%  {Phys. Rev. Lett.} {102}
%  ({2009}) {100502}. 
%  %,\urlprefix\url{http://link.aps.org/doi/10.1103/PhysRevLett.102.100502}.
%
%\bibitem[{Peschel and Eisler(2009)}]{Peschel2009}
%{I.~Peschel}, {V.~Eisler},
%  \textit{Reduced density matrices and entanglement entropy in free
%  lattice models}, {J. of Phys. A: Math. Theor.}
%  {42} ({2009}) {504003}.
%  %, \urlprefix\url{http://stacks.iop.org/1751-8121/42/i=50/a=504003}.
%
%\bibitem[{Turner et~al.(2010)Turner, Zhang, and Vishwanath}]{Turner2010}
%{M.~Turner, A}, {Y.~Zhang},
%  {A.~Vishwanath}, \textit{Entanglement and inversion
%  symmetry in topological insulators}, {Phys. Rev. B}
%  {82} ({2010}) {241102}.
%  %, \urlprefix\url{http://link.aps.org/doi/10.1103/PhysRevB.82.241102}.
%
%\bibitem[{Pastur and Slavin(2014)}]{Pastur2014}
%{L.~Pastur}, {V.~Slavin}, \textit{Area
%  Law Scaling for the Entropy of Disordered Quasifree Fermions},
%  {Phys. Rev. Lett.} {113}
%  ({2014}) {150404}. 
%  %,  \urlprefix\url{http://link.aps.org/doi/10.1103/PhysRevLett.113.150404}.
%
%\bibitem[{Potter(2014)}]{Potter2014}
%{A.~C. Potter}, \textit{{Boundary-law scaling of
%  entanglement entropy in diffusive metals}}, {arXiv:1408.1094},
%  {2014}.
%
%\bibitem[{Pouranvari et~al.(2015)Pouranvari, Zhang, and Yang}]{Pouranvari2015}
%{M.~Pouranvari}, {Y.~Zhang},
%  {K.~Yang}, \textit{Entanglement area law in
%  disordered free fermion Anderson model in one, two, and three dimensions},
%  {Adv. Condens. Matter Phys.} {2015}
%  ({2015}) {397630}.
%
%\bibitem[{Kopp et~al.(2007)Kopp, Jia, and Chakravarty}]{Kopp2007}
%{A.~Kopp}, {X.~Jia},
%  {S.~Chakravarty}, \textit{Replacing energy by von
%  {Neumann} entropy in quantum phase transitions}, {Ann. Phys.
%  (N.Y.)} {322} ({2007}) {1466}.
%  %, \urlprefix\url{http://www.sciencedirect.com/science/article/pii/S0003491606001825}.
%
%\bibitem[{Jia et~al.(2008)Jia, Subramaniam, Gruzberg, and
%  Chakravarty}]{Gruzberg2008}
%{X.~Jia}, {A.~R. Subramaniam},
%  {I.~A. Gruzberg}, {S.~Chakravarty},
%  \textit{Entanglement entropy and multifractality at localization
%  transitions}, {Phys. Rev. B} {77}
%  ({2008}) {014208}.
%  %, \urlprefix\url{http://link.aps.org/doi/10.1103/PhysRevB.77.014208}.
%
%\bibitem[{Berkovits(2012)}]{Berkovits2012}
%{R.~Berkovits}, \textit{Entanglement Entropy in a
%  One-Dimensional Disordered Interacting System: {The} Role of Localization},
%  {Phys. Rev. Lett.} {108}
%  ({2012}) {176803}.
%  %, \urlprefix\url{http://link.aps.org/doi/10.1103/PhysRevLett.108.176803}.
%
%\bibitem[{Zhao et~al.(2013)Zhao, Chu, and Shen}]{Zhao2013}
%{A.~Zhao}, {R.-L. Chu}, {S.-Q.
%  Shen}, \textit{Finite-size scaling of entanglement entropy at the
%  {Anderson} transition with interactions}, {Phys. Rev. B}
%  {87} ({2013}) {205140}.
%  %, \urlprefix\url{http://link.aps.org/doi/10.1103/PhysRevB.87.205140}.
%
%\bibitem[{Berkovits(2015)}]{Berkovits2015}
%{R.~Berkovits}, \textit{Entanglement Properties and
%  Quantum Phases for a Fermionic Disordered One-Dimensional Wire with
%  Attractive Interactions}, {Phys. Rev. Lett.}
%  {115} ({2015}) {206401}.
%  %, \urlprefix\url{http://link.aps.org/doi/10.1103/PhysRevLett.115.206401}.
%
%\bibitem[{Levitov and Shytov(2003)}]{LS-book}
%{L.~S. Levitov}, {A.~V. Shytov},
%  \textit{Green's functions. Theory and practice},
%  \bibinfo{publisher}{FizMatLit-Nauka}, \bibinfo{address}{Moscow},
%  \bibinfo{note}{in Russian}, {2003}.
%
%\bibitem[{Gioev(2006)}]{Gioev2006}
%{D.~Gioev}, \textit{Szeg{\"o} limit theorem for
%  operators with discontinuous symbols and applications to entanglement
%  entropy}, {Int. Math. Res. Notices} {2006}.
%  %, \urlprefix\url{http://imrn.oxfordjournals.org/content/2006/O95181.abstract}.
%
%\bibitem[{Lee and Ramakrishnan(1985)}]{Lee1985}
%{P.~A. Lee}, {T.~V. Ramakrishnan},
%  \textit{Disordered electronic systems}, {Rev. Mod.
%  Phys.} {57} ({1985}) {287}.
%
%\bibitem[{Zala et~al.(2001)Zala, Narozhny, and Aleiner}]{Aleiner2001}
%{G.~Zala}, {B.~N. Narozhny},
%  {I.~L. Aleiner}, \textit{Interaction corrections at
%  intermediate temperatures: Longitudinal conductivity and kinetic equation},
%  {Phys. Rev. B} {64} ({2001})
%  {214204}.
%  %, \urlprefix\url{http://link.aps.org/doi/10.1103/PhysRevB.64.214204}.
%
%\bibitem[{Wegner(1976)}]{Wegner1976}
%{F.~Wegner}, \textit{Electrons in Disordered Systems.
%  {Scaling} near the Mobility Edge}, {Z. Phys. B}
%  {25} ({1976}) {327}.
%
%\bibitem[{Mott(1968)}]{Mott1968}
%{N.~F. Mott}, \textit{Conduction in non-crystalline
%  systems. {I. Localized} electronic states in disordered systems},
%  {Philos. Mag.} {17} (1968) 1259.
%  %, \urlprefix\url{http://www.tandfonline.com/doi/abs/10.1080/14786436808223200?journalCode=tphm19}.
%
%\bibitem[{Mott(1970)}]{Mott1970}
%{N.~F. Mott}, \textit{Conduction in non-crystalline
%  systems. {IV. Anderson} localization in a disordered lattice.},
%  {Philos. Mag.} {22} (1970) 7. 
%  %, \urlprefix\url{http://www.tandfonline.com/doi/abs/10.1080/14786437008228147}.
%
%\bibitem[{Burmistrov et~al.(2014)Burmistrov, Gornyi, and
%  Mirlin}]{Burmistrov2014}
%{I.~S. Burmistrov}, {I.~V. Gornyi},
%  {A.~D. Mirlin}, \textit{Tunneling into the localized
%  phase near {Anderson} transitions with {Coulomb} interaction},
%  {Phys. Rev. B} {89} ({2014})
%  {035430}.
%  %, \urlprefix\url{http://link.aps.org/doi/10.1103/PhysRevB.89.035430}.
%
%\bibitem[{Levitov and Lesovik(1993)}]{Levitov1993}
%{L.~S. Levitov}, {G.~B. Lesovik},
%  \textit{Charge distribution in quantum shot noise},
%  {JETP Lett.} {58} ({1993})
%  {230}.
%  %, \urlprefix\url{http://jetpletters.ac.ru/ps/1186/article_17907.shtml}.
%
%\bibitem[{Ivanov et~al.(2013)Ivanov, Abanov, and Cheianov}]{Ivanov2013}
%{D.~Ivanov}, {A.~Abanov},
%  {V.~Cheianov}, \textit{Counting free fermions on a
%  line: a Fisher–Hartwig asymptotic expansion for the Toeplitz determinant in
%  the double-scaling limit}, {J. Phys. A: Math. Theor}
%  {46} ({2013}) {085003}.
%  %, \urlprefix\url{http://iopscience.iop.org/article/10.1088/1751-8113/46/8/085003/meta}.
%
%\bibitem[{Song et~al.(2011)Song, Flindt, Rachel, Klich, and
%  Le~Hur}]{Flindt2011}
%{H.~F. Song}, {C.~Flindt},
%  {S.~Rachel}, {I.~Klich},
%  {K.~Le~Hur}, \textit{Entanglement entropy from charge
%  statistics: {Exact} relations for noninteracting many-body systems},
%  {Phys. Rev. B} {83} ({2011})
%  {161408}.
%  %,  \urlprefix\url{http://link.aps.org/doi/10.1103/PhysRevB.83.161408}.
%
%\bibitem[{Calabrese et~al.(2012)Calabrese, Mintchev, and
%  Vicari}]{Calabrese2012}
%{P.~Calabrese}, {M.~Mintchev},
%  {E.~Vicari}, \textit{Exact relations between particle
%  fluctuations and entanglement in Fermi gases}, {EPL
%  (Europhysics Letters)} {98} ({2012})
%  {20003}.
%  %, \urlprefix\url{http://stacks.iop.org/0295-5075/98/i=2/a=20003}.
%
%\bibitem[{Song et~al.(2012)Song, Rachel, Flindt, Klich, Laflorencie, and
%  Le~Hur}]{Flindt2012}
%{H.~F. Song}, {S.~Rachel},
%  {C.~Flindt}, {I.~Klich},
%  {N.~Laflorencie}, {K.~Le~Hur},
%  \textit{Bipartite fluctuations as a probe of many-body entanglement},
%  {Phys. Rev. B} {85} ({2012})
%  {035409}.
%  %,  \urlprefix\url{http://link.aps.org/doi/10.1103/PhysRevB.85.035409}.
%
%\bibitem[{Thomas and Flindt(2015)}]{Flindt2015}
%{K.~H. Thomas}, {C.~Flindt},
%  \textit{Entanglement entropy in dynamic quantum-coherent conductors},
%  {Phys. Rev. B} {91} ({2015})
%  {125406}.
%  %, \urlprefix\url{http://link.aps.org/doi/10.1103/PhysRevB.91.125406}.
%
%\bibitem[{Zhao and Bhatt(2016)}]{Bhatt2016}
%{L.~Zhao}, {R.~N. Bhatt},
%  \textit{Quantum Entanglement as a Diagnostic of Phase Transitions in
%  Disordered Fractional Quantum {Hall} Liquids},
% Phys. Rev. Lett. 117 (2016) 206801.
%
%\end{thebibliography}

\end{document}


